\documentclass[11pt]{article}% Shared preamble for the rigor documents (Lamport-style structured proofs).  \input this file after \documentclass.
\usepackage[T1]{fontenc}\usepackage{lmodern}\usepackage{microtype}
\usepackage[margin=1in]{geometry}
\usepackage{amsmath,amssymb,amsthm,mathtools}
\usepackage{xcolor}\usepackage{enumitem}\usepackage{booktabs}\usepackage{graphicx}\usepackage{hyperref}
\usepackage[most]{tcolorbox}
\definecolor{navy}{HTML}{17365D}\definecolor{teal}{HTML}{117A8B}\definecolor{warn}{HTML}{A33A2B}\definecolor{okgreen}{HTML}{2E7D4F}
\hypersetup{colorlinks=true,linkcolor=navy,citecolor=teal,urlcolor=teal}
\theoremstyle{plain}
\newtheorem{theorem}{Theorem}[section]\newtheorem{lemma}[theorem]{Lemma}\newtheorem{proposition}[theorem]{Proposition}\newtheorem{corollary}[theorem]{Corollary}
\theoremstyle{definition}\newtheorem{definition}[theorem]{Definition}\newtheorem{remark}[theorem]{Remark}\newtheorem{claim}[theorem]{Claim}\newtheorem{issue}[theorem]{Issue}
% ---------- Lamport structured proofs ----------
% \lstep{level}{number}{assertion}   -> indented "<level>number. assertion"
% \lproof                             -> "PROOF:" line (start of the sub-proof of the preceding step)
% \lby{justification}                 -> "BY ..." leaf justification for the preceding step
% \lqed{level}{number}                -> QED step of a sub-proof
% keywords: \lassume \lprove \lsuffices \lcase \ldefine \lpick \lnew
\newlength{\lindent}\setlength{\lindent}{1.7em}
\newcommand{\lstep}[3]{\par\smallskip\noindent\hspace*{\dimexpr\numexpr#1-1\relax\lindent\relax}{\color{navy}$\langle#1\rangle#2$.}\ #3}
\newcommand{\lproof}{\par\noindent\hspace*{\lindent}{\scshape Proof:}}
\newcommand{\lby}[1]{\par\noindent\hspace*{\lindent}{\scshape By}\ #1.}
\newcommand{\lqed}[2]{\par\smallskip\noindent\hspace*{\dimexpr\numexpr#1-1\relax\lindent\relax}{\color{navy}$\langle#1\rangle#2$.}\ {\scshape Q.E.D.}}
\newcommand{\lassume}{{\scshape Assume}\ }\newcommand{\lprove}{{\scshape Prove}\ }\newcommand{\lsuffices}{{\scshape Suffices}\ }
\newcommand{\lcase}{{\scshape Case}\ }\newcommand{\ldefine}{{\scshape Define}\ }\newcommand{\lpick}{{\scshape Pick}\ }\newcommand{\lnew}{{\scshape New}\ }
% ---------- verdict boxes ----------
\newtcolorbox{verdictok}{colback=okgreen!8,colframe=okgreen,title={\bfseries Verdict: verified},fonttitle=\small}
\newtcolorbox{verdictgap}{colback=orange!10,colframe=orange!80!black,title={\bfseries Verdict: gap / caveat},fonttitle=\small}
\newtcolorbox{verdictbreak}{colback=warn!10,colframe=warn,title={\bfseries Verdict: proof-breaking issue},fonttitle=\small}
\newtcolorbox{citedbox}[1]{colback=navy!5,colframe=navy!60,title={\bfseries Cited result: #1},fonttitle=\small,breakable}
\setlength{\parindent}{0pt}\setlength{\parskip}{4pt}\allowdisplaybreaks
\newcommand{\cA}{\mathcal A}\newcommand{\cF}{\mathcal F}\newcommand{\cM}{\mathfrak M}\newcommand{\cN}{\mathfrak N}

% extra calligraphic shortcuts (\providecommand: no clash if a document defines its own)
\providecommand{\cB}{\mathcal B}\providecommand{\cC}{\mathcal C}\providecommand{\cD}{\mathcal D}\providecommand{\cE}{\mathcal E}\providecommand{\cG}{\mathcal G}\providecommand{\cH}{\mathcal H}\providecommand{\cI}{\mathcal I}\providecommand{\cJ}{\mathcal J}\providecommand{\cK}{\mathcal K}\providecommand{\cL}{\mathcal L}\providecommand{\cO}{\mathcal O}\providecommand{\cP}{\mathcal P}\providecommand{\cQ}{\mathcal Q}\providecommand{\cR}{\mathcal R}\providecommand{\cS}{\mathcal S}\providecommand{\cT}{\mathcal T}\providecommand{\cU}{\mathcal U}\providecommand{\cV}{\mathcal V}\providecommand{\cW}{\mathcal W}\providecommand{\cX}{\mathcal X}\providecommand{\cY}{\mathcal Y}\providecommand{\cZ}{\mathcal Z}
\newcommand{\Fid}{\operatorname{F}}\newcommand{\Tr}{\operatorname{Tr}}\newcommand{\eps}{\varepsilon}\newcommand{\dd}{\,\mathrm d}

\newcommand{\Wh}{\widehat W}\newcommand{\gh}{\widehat{\widetilde g}}\newcommand{\gt}{\widetilde g}\newcommand{\Diff}{\mathrm{Diff}}
\title{\vspace{-1.4cm}\color{navy}\bfseries Referee report V5:\\ Note~7 (\texttt{universality\_normalization}) and the mathematical claims of Note~6 (\texttt{strategy\_open\_problems})\\[.3em]\large Lamport-style verification of $f_2=c/(12\pi^2)$, $s_2=c/12$}
\author{}\date{2026-09-05 (text as of the 16:13 errata)}
\begin{document}\maketitle
\vspace{-1.2em}
\begin{abstract}\noindent
Every displayed formula of Note~7 that enters the derivation of $f_2=c/(12\pi^2)$ and $s_2=c/12$ has been
re-derived from scratch and, where possible, independently checked numerically.  \textbf{The two closed forms are
correct.}  The load-bearing normalization $\langle T(x)T(y)\rangle=\frac{c}{8\pi^2}(x-y-i0)^{-4}$ is confirmed by an
explicit Wick computation for the free complex chiral fermion in the generator normalization (Theorem~\ref{thm:TT}
below), independently of the note; so are $\widehat W$, $\widehat{\widetilde g}$, the kernel of the first-order
defect, the two spectral integrals, and the whole factor bookkeeping of Theorem~A.  The auxiliary integral
$\int w^2(\cosh w+\cosh a)^{-1}dw=\frac{2a(a^2+\pi^2)}{3\sinh a}$, which the note only verifies numerically, is
proved here by residues.  What remains open is not arithmetic but functional analysis: for the zero-collar
representative $T(g)\Omega\notin\mathcal H$ and $T(g)\notin\cM$, so Lemma~3.1(2) --- whose KMS symmetry Sections~4.2--4.3
use --- is not applicable to it, and the general route of Sections~3--4 is a formal computation validated
\emph{a posteriori} by the free-fermion route of Section~5.  Three further defects are recorded: a false
boundedness claim after eq.~(13); a modular-flow sign left as ``$\pm$'' and then fixed by an argument that
presupposes what is at issue (repairable --- the sign is derived here from Bisognano--Wichmann); and, in Note~6,
an internal clash between the BPZ two-point function in \S2 step~3 and the generator convention in step~4,
plus an energy formula that is the energy of $U(\varphi^{-1})\Omega$, not of $U(\varphi)\Omega$.
\end{abstract}
\tableofcontents
\bigskip
\begin{center}\footnotesize
\begin{tabular}{llp{6.3cm}}\toprule
Location & Verdict & One line\\\midrule
N7 eq.~(1) $\langle TT\rangle=\frac{c}{8\pi^2}u^{-4}$ & \textbf{verified} & independent Wick computation, $c=1$ Dirac\\
N7 eq.~(2),(3) $\xi$-frame, $W(u)$ & \textbf{verified} & incl.\ $(x-y)/\sqrt{\beta\beta'}=2\sinh\frac{\xi-\xi'}2$\\
N7 \S1 ``$\xi\mapsto\xi\pm2\pi t$'' & gap & sign left open; it is $\xi\mapsto\xi-2\pi t$\\
N7 Lemma~2.1, eq.~(6),(7) & \textbf{verified} & $\widetilde g=-4\sinh^2\frac\xi2$, $\widehat{\widetilde g}=\frac{-2i}{k(k^2+1)}$\\
N7 eq.~(8),(9) defect kernel & \textbf{verified} & sign of $\mathfrak g$ and $\widetilde D^{(1)}$ re-derived\\
N7 Lemma~3.1 (1),(2),(3) & \textbf{verified} & all standard-form identities check out\\
N7 text after eq.~(13) & \textbf{break (local)} & ``both [weights] are bounded by $1+\lambda$'' is false for KM\\
N7 Lemma~3.2 & gap & one hypothesis vacuous; not applicable to the note's own $A$\\
N7 Lemma~4.1, eq.~(14) & \textbf{verified} & contour shift legitimate; $\widehat W$ correct\\
N7 \S4.2 eq.~(15) & gap & KMS sign-fixing presupposes $T(g)\in\cM$, which is false\\
N7 \S4.3 eq.~(16)--(18) & \textbf{verified} & all prefactors and the sum $=2$\\
N7 \S4.4 universality & gap & logic valid but the load is carried by \S5, not by ``agreement''\\
N7 Lemma~5.1, eq.~(20)--(23) & \textbf{verified} & pointwise identity, weights, $|\mathfrak g|^2$, $V(u)$\\
N7 Theorem~5.2 & \textbf{verified} & assembly correct, $4/\pi^2\to1/\pi^2$ erratum confirmed\\
N7 Cor.~6.1 & \textbf{verified} & VWZ conversion and $12/(\pi^2c)$ correct (chiral)\\
N6 \S2 step~3 ($\frac c{32}$) & \textbf{break (local)} & BPZ constant inside a generator-normalized argument\\
N6 \S2 eq.~(2),(4) & \textbf{verified} & metrics and both integrals correct\\
N6 \S3 eq.~(5) energy & minor & it is $E(\varphi^{-1})$; combination does vanish on $\mathrm{PSL}(2,\mathbb R)$\\
N6 \S3 Route~B & gap & unbounded ``compact'' set; two uncited theorems; strong additivity\\
N6 \S4.1 $c/18$ per chirality & \textbf{verified} & consistent with the two-branch law\\
\bottomrule\end{tabular}\end{center}
\clearpage
\section{Section 1 of Note 7: conventions and the two-point normalization}\label{sec:conv}

The whole numerical value of $f_2$ and $s_2$ hangs on the constant in
\begin{quote}\itshape
``$T(f)=\int T(x)f(x)\,dx$ is normalized as the \emph{generator}: $e^{isT(f)}$ implements the flow of the vector
field $f\partial_x$, so that $T(1)$ is the light-ray momentum.  In this normalization
$\langle\Omega,T(x)T(y)\Omega\rangle=\frac{c}{8\pi^2}\frac1{(x-y-i0)^4}$'' \hfill (N7, eq.~(1)).
\end{quote}
A factor $\alpha$ here multiplies $\widehat W$, hence $g_B$ and $g_{\mathrm{KM}}$, hence $f_2$ and $s_2$, by
$\alpha$.  We verify it from first principles for the free complex (Dirac) chiral fermion, for which $c=1$.

\begin{definition}[Fermion conventions]\label{def:ferm}
$\psi$ is the chiral Fermi field on $\mathbb R$ with $\{\psi(x),\psi^\dagger(y)\}=\delta(x-y)$,
$\{\psi(x),\psi(y)\}=0$, acting on the Fock space of the Fermi sea, $\Omega$ the vacuum.  Write
$\psi(x)=\int\frac{dp}{2\pi}e^{ipx}c(p)$ with $\{c(p),c(q)^\dagger\}=2\pi\delta(p-q)$.  The stress tensor is
$T=\tfrac i2\bigl(:\!\psi^\dagger\partial\psi\!:-:\!\partial\psi^\dagger\,\psi\!:\bigr)$.
\end{definition}

\begin{theorem}[Normalization of eq.~(1)]\label{thm:TT}
With Definition~\ref{def:ferm} and the requirement that $P:=T(\mathbf 1)$ satisfy $[P,\psi(x)]=-i\psi'(x)$,
$P\ge0$, $P\Omega=0$:
\begin{equation}\label{eq:v-prop}
 \langle\Omega,\psi(x)\psi^\dagger(y)\Omega\rangle=\langle\Omega,\psi^\dagger(x)\psi(y)\Omega\rangle
 =\frac1{2\pi i}\,\frac1{x-y-i0}=:A(x-y),
\end{equation}
\begin{equation}\label{eq:v-TT}
 \langle\Omega,T(x)T(y)\Omega\rangle=\tfrac12\bigl[A'(u)^2-A(u)A''(u)\bigr]=\frac1{8\pi^2}\,\frac1{(x-y-i0)^4},
 \qquad u=x-y .
\end{equation}
Hence N7 eq.~(1) holds with $c=1$ for the Dirac fermion, i.e.\ with $c$ the central charge in the convention
$c=r$ for $r$ complex components.
\end{theorem}

\lstep{1}{1}{$[P,c(p)]=p\,c(p)$, and therefore $c(p)\Omega=0$ for $p<0$ and $c(p)^\dagger\Omega=0$ for $p>0$.}
\lproof
\lstep{2}{1}{$-i\psi'(x)=-i\int\frac{dp}{2\pi}(ip)e^{ipx}c(p)=\int\frac{dp}{2\pi}p\,e^{ipx}c(p)$.}
\lby{differentiating under the integral sign in Definition~\ref{def:ferm}}
\lstep{2}{2}{Comparing $\langle2\rangle1$ with $[P,\psi(x)]=\int\frac{dp}{2\pi}e^{ipx}[P,c(p)]$ and using that the
$e^{ipx}$ are a complete system gives $[P,c(p)]=p\,c(p)$.}
\lby{$\langle2\rangle1$ and Fourier inversion}
\lstep{2}{3}{$c(p)$ raises the $P$-eigenvalue by $p$; since $P\ge0$ and $P\Omega=0$, $c(p)\Omega$ has
$P$-eigenvalue $p$, which is inadmissible for $p<0$; likewise $c(p)^\dagger$ lowers it by $p$, inadmissible for
$p>0$.}
\lby{$\langle2\rangle2$, $P\ge0$, $P\Omega=0$}
\lqed{2}{4}
\lby{$\langle2\rangle3$}
\lstep{1}{2}{$\langle\psi(x)\psi^\dagger(y)\rangle=\int_{-\infty}^0\frac{dp}{2\pi}e^{ip(x-y)}=\frac{1}{2\pi i}(x-y-i0)^{-1}$.}
\lproof
\lstep{2}{1}{$\langle c(p)c(q)^\dagger\rangle=2\pi\delta(p-q)\mathbf 1_{q<0}$.}
\lby{$\langle1\rangle1$ and the CAR of Definition~\ref{def:ferm}}
\lstep{2}{2}{$\int_{-\infty}^0e^{ipu+\eps p}dp=\frac1{iu+\eps}=\frac{-i}{u-i\eps}$ for $\eps>0$.}
\lby{elementary integration; the regulator $e^{\eps p}$ is the adiabatic cutoff, and the limit $\eps\downarrow0$
is taken in $\mathcal S'(\mathbb R)$}
\lqed{2}{3}
\lby{$\langle2\rangle1$, $\langle2\rangle2$}
\lstep{1}{3}{$\langle\psi^\dagger(x)\psi(y)\rangle=A(x-y)$ as well.}
\lproof
\lstep{2}{1}{$\langle\psi^\dagger(y)\psi(x)\rangle=\int_0^\infty\frac{dp}{2\pi}e^{ip(x-y)}=\frac{i}{2\pi}(x-y+i0)^{-1}$.}
\lby{$\langle1\rangle1$ (now $c(p)\Omega=0$ for $p<0$) and $\int_0^\infty e^{ipu-\eps p}dp=\frac i{u+i\eps}$}
\lstep{2}{2}{Exchanging the names $x\leftrightarrow y$ in $\langle2\rangle1$:
$\langle\psi^\dagger(x)\psi(y)\rangle=\frac i{2\pi}(y-x+i0)^{-1}=\frac{-i}{2\pi}(x-y-i0)^{-1}=A(x-y)$.}
\lby{$\langle2\rangle1$}
\lqed{2}{3}
\lby{$\langle2\rangle2$}
\lstep{1}{4}{The conventions are consistent: $\langle\{\psi(x),\psi^\dagger(y)\}\rangle
=\frac1{2\pi i}\bigl[(x-y-i0)^{-1}-(x-y+i0)^{-1}\bigr]=\delta(x-y)$.}
\lby{$\langle1\rangle2$, $\langle1\rangle3$ and the Sokhotski formula $(u\mp i0)^{-1}=\mathrm{PV}\frac1u\pm i\pi\delta(u)$}
\lstep{1}{5}{$T(\mathbf 1)=i\int:\!\psi^\dagger\psi'\!:dx$ and $[T(\mathbf 1),\psi(y)]=-i\psi'(y)$, so the
normalization of Definition~\ref{def:ferm} is the generator normalization of N7 \S1.}
\lproof
\lstep{2}{1}{$\int(-\partial\psi^\dagger\,\psi)dx=\int\psi^\dagger\partial\psi\,dx$, hence
$T(\mathbf 1)=\frac i2\int(\psi^\dagger\psi'-\psi^{\dagger\prime}\psi)=i\int\psi^\dagger\psi'$.}
\lby{integration by parts; the boundary terms vanish in the sense of distributions on test functions of compact
support, and for $f\equiv\mathbf 1$ on the Fock domain of finite-energy vectors}
\lstep{2}{2}{$[\psi^\dagger(x)\psi'(x),\psi(y)]=\psi^\dagger(x)\{\psi'(x),\psi(y)\}-\{\psi^\dagger(x),\psi(y)\}\psi'(x)=-\delta(x-y)\psi'(x)$.}
\lby{$[AB,C]=A\{B,C\}-\{A,C\}B$ and the CAR}
\lqed{2}{3}
\lby{$\langle2\rangle1$, $\langle2\rangle2$: $[T(\mathbf 1),\psi(y)]=i\int(-\delta(x-y)\psi'(x))dx=-i\psi'(y)$}
\lstep{1}{6}{$\langle\Omega,T(x)T(y)\Omega\rangle=\frac12\bigl[A'(u)^2-A(u)A''(u)\bigr]$.}
\lproof
\lstep{2}{1}{Write $T=\frac i2\sum_a s_a:\!\partial^{m_a}\psi^\dagger\,\partial^{n_a}\psi\!:$ with
$(m,n,s)=(0,1,+1)$ and $(1,0,-1)$.}
\lby{Definition~\ref{def:ferm}}
\lstep{2}{2}{For each pair $(a,b)$ the connected Wick contraction is
$\langle\partial^{m_a}\psi^\dagger(x)\,\partial^{n_b}\psi(y)\rangle\,\langle\partial^{n_a}\psi(x)\,\partial^{m_b}\psi^\dagger(y)\rangle
=(-1)^{m_b+n_b}A^{(m_a+n_b)}(u)A^{(n_a+m_b)}(u)$, with a $+$ sign.}
\lproof
\lstep{3}{1}{$\langle 1234\rangle=\langle12\rangle\langle34\rangle-\langle13\rangle\langle24\rangle+\langle14\rangle\langle23\rangle$
with $1=\psi^\dagger(x),2=\psi(x),3=\psi^\dagger(y),4=\psi(y)$; the middle term vanishes because
$\langle\psi^\dagger\psi^\dagger\rangle=\langle\psi\psi\rangle=0$, and the first is the disconnected term removed
by normal ordering.}
\lby{Wick's theorem for quasi-free states and $\langle\psi\psi\rangle=0$ in the gauge-invariant vacuum}
\lstep{3}{2}{Differentiating $\langle\psi^\dagger(x)\psi(y)\rangle=A(x-y)$ (step $\langle1\rangle3$) gives
$\partial_x^m\partial_y^{n}A(x-y)=(-1)^{n}A^{(m+n)}(u)$; same for $\langle\psi(x)\psi^\dagger(y)\rangle$.}
\lby{$\langle1\rangle2$, $\langle1\rangle3$, chain rule}
\lqed{3}{3}\lby{$\langle3\rangle1$, $\langle3\rangle2$}
\lstep{2}{3}{$m_b+n_b=1$ for both $b$, so $(-1)^{m_b+n_b}=-1$; with the prefactor $(\frac i2)^2=-\frac14$ one gets
$\langle TT\rangle_c=\frac14\sum_{a,b}s_as_b\,A^{(m_a+n_b)}A^{(n_a+m_b)}$.}
\lby{$\langle2\rangle1$, $\langle2\rangle2$}
\lstep{2}{4}{The four terms are $(A')^2,\;-AA'',\;-A''A,\;(A')^2$, so
$\langle TT\rangle_c=\frac14\bigl[2(A')^2-2AA''\bigr]$.}
\lby{$\langle2\rangle3$ with $(m_a,n_a),(m_b,n_b)\in\{(0,1),(1,0)\}$}
\lqed{2}{5}\lby{$\langle2\rangle4$}
\lstep{1}{7}{With $A=a/u$, $a=(2\pi i)^{-1}$: $(A')^2-AA''=\frac{a^2}{u^4}-\frac{2a^2}{u^4}=-\frac{a^2}{u^4}$, so
$\langle TT\rangle=-\frac{a^2}{2u^4}=\frac1{8\pi^2}\frac1{u^4}$, $u=x-y-i0$.}
\lby{$\langle1\rangle6$, $A'=-a/u^2$, $A''=2a/u^3$, $a^2=-1/(4\pi^2)$}
\lqed{1}{8}
\lby{$\langle1\rangle5$, $\langle1\rangle7$}

\begin{verdictok}
N7 eq.~(1) is correct.  Numerically: $-a^2/2=0.0126651479552922=1/(8\pi^2)$ to 15 digits, and a finite-difference
evaluation of $\frac12[(A')^2-AA'']$ at $u=1.3-0.001i$ reproduces $\frac1{8\pi^2}u^{-4}$ to 10 digits
(\texttt{numerics/v5chk.py}).  Two independent corroborations follow below: the flat-space Fourier transform
in Lemma~\ref{lem:What7} and, decisively, the free-fermion computation of N7 \S5, which produces
$g_B=\frac2{3\pi^2}$ with no free constant and therefore \emph{fixes} $\alpha=\frac1{8\pi^2}$ at $c=1$.
The parenthesis ``the BPZ field $2\pi T$ has $\frac c2(x-y)^{-4}$'' is correct: $(2\pi)^2\frac c{8\pi^2}=\frac c2$.
\end{verdictok}

\begin{verdictok}
$T(f)\Omega=0$ for $f\in\operatorname{span}\{1,x,x^2\}$ is correct: these are exactly the generators of the
M\"obius subgroup $\mathrm{PSL}(2,\mathbb R)$ of $\mathrm{Diff}_+(S^1)$ acting on the line (translation, dilation, special
conformal), $U(s)\Omega=\Omega$ for the corresponding one-parameter groups, and $T(f)=-i\frac{d}{ds}U(s)|_{s=0}$
on the domain of the generator.  Note that this is the only place where the \emph{unsubtracted} $T$ occurs;
$\langle T(x)\rangle=0$ on the line is consistent with $T(\mathbf 1)\Omega=P\Omega=0$.
\end{verdictok}

\begin{lemma}[N7 eq.~(2)--(3)]\label{lem:W}
With $\xi=-\log\frac{L-x}{L}$, $\beta=\frac{dx}{d\xi}=L-x=Le^{-\xi}$, $x=L(1-e^{-\xi})$ one has
$\frac{x-y}{\sqrt{\beta\beta'}}=2\sinh\frac{\xi-\xi'}2$ and, with $\widetilde T(\xi)=\beta^2T(x)$,
$W(u)=\langle\widetilde T\widetilde T\rangle_{\mathrm{conn}}=\frac{c}{128\pi^2}\sinh^{-4}\frac{u-i0}2$.
\end{lemma}
\lstep{1}{1}{$\beta=dx/d\xi=Le^{-\xi}=L-x$ and $\xi<0\Leftrightarrow x<0$, $\xi>0\Leftrightarrow 0<x<L$.}
\lby{differentiating $x=L(1-e^{-\xi})$; $\xi\mapsto x$ is an increasing bijection $\mathbb R\to(-\infty,L)$}
\lstep{1}{2}{$x-y=Le^{-(\xi+\xi')/2}\bigl(e^{(\xi-\xi')/2}-e^{-(\xi-\xi')/2}\bigr)=2\sqrt{\beta\beta'}\sinh\frac{\xi-\xi'}2$.}
\lby{$x-y=L(e^{-\xi'}-e^{-\xi})$ and $\sqrt{\beta\beta'}=Le^{-(\xi+\xi')/2}$}
\lstep{1}{3}{$\widetilde T(\xi)=\beta^2T(x)$ is the weight-2 density: $T(f)=\int T(x)f(x)dx=\int\widetilde T(\xi)\widetilde g(\xi)d\xi$
with $\widetilde g=g/\beta$.}
\lby{$T(x)g(x)dx=\beta^{-2}\widetilde T(\xi)g(x(\xi))\beta\,d\xi$}
\lstep{1}{4}{$\dfrac{\beta^2\beta'^2}{(x-y-i0)^4}=\Bigl(\dfrac{\sqrt{\beta\beta'}}{x-y-i0}\Bigr)^4
=\dfrac1{16\sinh^4\frac{u-i0}2}$, so $W(u)=\frac{c}{8\pi^2}\cdot\frac1{16}\sinh^{-4}\frac{u-i0}2$.}
\lby{$\langle1\rangle2$, eq.~(1); $\beta\beta'>0$, so the $-i0$ prescription passes unchanged to $u$}
\lqed{1}{5}\lby{$\langle1\rangle4$, $\frac{c}{128\pi^2}=\frac{c}{8\pi^2\cdot16}$}
\begin{verdictok}
Eqs.~(2) and (3) verified, including the $i0$ prescription.  $\langle\widetilde T\rangle=\beta^2\langle T\rangle=0$,
so the subtraction in ``$\widetilde T_0=\widetilde T-\langle\widetilde T\rangle$'' (N7, after eq.~(6)) is vacuous
in the note's own convention $\widetilde T=\beta^2T$; it would be non-vacuous only if $\widetilde T$ were defined
with the Schwarzian term.  Harmless either way: the difference is a $c$-number, which drops out of the commutator
of Lemma~2.1 and, since $(S^*-1)\Omega=0$, also out of $\eta$.
\end{verdictok}
\subsection{The modular frame: the sign of the flow}\label{sec:bwsign}
N7 \S1 says: ``\emph{The modular group of $(\cA(I),\Omega)$ is the dilation about $L$ (Bisognano--Wichmann), a
translation $\xi\mapsto\xi\pm2\pi t$ for $\Delta^{it}$}''.  The $\pm$ is not innocuous: it decides whether the
spectral parameter is $u=+2\pi k$ or $u=-2\pi k$ in eq.~(15), i.e.\ whether the Bures weight is evaluated at
$\lambda=e^{2\pi k}$ or $e^{-2\pi k}$.  N7 \S4.2 fixes it by KMS; \S\ref{sec:kmssign} below shows that this
argument presupposes $T(g)\in\cM$, which is false.  The sign is however determined outright:

\begin{lemma}[Modular flow of the half-line $I=(-\infty,L)$]\label{lem:bw}
For a M\"obius covariant net on the line with Haag duality, $\Delta_{\cA(I)}^{it}$ implements the point map
$\varphi_t(x)=L+e^{2\pi t}(x-L)$, i.e.\ $\xi\mapsto\xi-2\pi t$, and consequently
$\Delta^{it}\widetilde T_0(\xi)\Delta^{-it}=\widetilde T_0(\xi-2\pi t)$.
\end{lemma}
\lstep{1}{1}{For $M=\cA((0,\infty))$, $\Delta_M^{it}$ implements $x\mapsto e^{-2\pi t}x$.}
\lby{Bisognano--Wichmann for conformal nets: Brunetti--Guido--Longo, CMP \textbf{156} (1993), Thm.~2.4 /
Prop.~1.1 (\texttt{refs/BGL93\_funct-an-9302008.pdf}); equivalently Borchers, CMP \textbf{143} (1992):
$\Delta^{it}U(a)\Delta^{-it}=U(e^{-2\pi t}a)$, which is exactly the conjugation of a translation by the
dilation $x\mapsto e^{-2\pi t}x$}
\lstep{1}{2}{$\cA((-\infty,L))=\cA((L,\infty))'$ and $\Delta_{M'}=\Delta_M^{-1}$.}
\lby{Haag duality on the line for M\"obius covariant nets (BGL93, Thm.~2.19(b)) and Tomita--Takesaki}
\lstep{1}{3}{Hence $\Delta_{\cA(I)}^{it}=\Delta_{\cA((L,\infty))}^{-it}$ implements $x-L\mapsto e^{+2\pi t}(x-L)$,
i.e.\ $L-x\mapsto e^{2\pi t}(L-x)$, i.e.\ $\beta\mapsto e^{2\pi t}\beta$, i.e.\ $\xi=-\log(\beta/L)\mapsto\xi-2\pi t$.}
\lby{$\langle1\rangle1$, $\langle1\rangle2$ and eq.~(2)}
\lstep{1}{4}{If $U$ implements the point map $\phi$ (so that $U\cA(I)U^*=\cA(\phi(I))$) then
$U\,O(x)\,U^*=\phi'(x)^{h}O(\phi(x))$ for a weight-$h$ density.}
\lby{fixed by two special cases: $U(a)=e^{iaP}$, $[P,\psi]=-i\psi'$ gives $U(a)\psi(x)U(a)^*=\psi(x+a)$;
the dilation $\phi(x)=\lambda x$ requires $U_\lambda\psi(x)U_\lambda^*=\lambda^{1/2}\psi(\lambda x)$ for the CAR
$\{\cdot,\cdot\}=\delta$ to be preserved, which is $h=\frac12$}
\lqed{1}{5}
\lby{$\langle1\rangle3$, $\langle1\rangle4$ with $h=2$, $\phi_t'\equiv1$ in the $\xi$-chart; the Schwarzian of the
translation vanishes}
\begin{verdictgap}
The $\pm$ should be a $-$.  The consequence is that
$\langle T(g)\Omega,\Delta^{it}T(g)\Omega\rangle=\iint\widetilde g\widetilde g'\,W(\xi-\xi'+2\pi t)
=\frac1{2\pi}\int dk\,\widehat W(k)|\widehat{\widetilde g}(k)|^2e^{+2\pi ikt}$, i.e.\ $u=\log\lambda=+2\pi k$,
which is the sign N7 uses.  \textbf{The note's answer is right; the argument it gives for it is not
(see \S\ref{sec:kmssign}), but this Lemma repairs it.}  Repair: replace ``$\pm$'' by ``$-$'' and cite BGL93
Thm.~2.19(b) for Haag duality; drop the appeal to KMS in \S4.2.
\end{verdictgap}

\section{Section 2 of Note 7: the tangent vector}\label{sec:tangent}

\begin{verdictgap}
\textbf{Lemma~2.1} is correct as an identity of functionals but its proof has two soft spots.
(i) ``\emph{both sides are normal in $X$ because $T(g)\Omega$ is a vector ($g$ is smooth and compactly
supported)}'': for the collar representative $g=f\chi$ this is fine, and normality of
$X\mapsto i\langle T(g)\Omega,X\Omega\rangle-i\langle S^*T(g)\Omega,X\Omega\rangle$ follows from
$T(g)\Omega\in\mathcal H$; but ``$g$ smooth and compactly supported $\Rightarrow$ $T(g)\Omega\in\mathcal H$''
needs the linear energy bound $\|T(f)\psi\|\le C\|f\|_{3/2}\|(1+L_0)\psi\|$ of Carpi--Weiner (= CDIT
Prop.~2.2(4)), not Fewster--Hollands (cf.\ finding [R4] on N6).  It is \emph{not} available for the zero-collar
representative used from eq.~(6) on --- see \S\ref{sec:normdiv}.
(ii) ``$\frac d{ds}\omega(x_AU(s)y_DU(s)^*)|_0=i\omega(x_A[T(f),y_D])$'' requires differentiability of
$s\mapsto U(s)y_DU(s)^*$ in the relevant topology; $U(s)$ is a M\"obius one-parameter group and $y_D\in\cA(D)$,
so this is standard, but $T(f)$ is unbounded and the commutator is to be read as a form on $\cA(D)\Omega$.
Neither point affects the conclusion.
\end{verdictgap}

\begin{verdictok}
\textbf{The second representative.}  ``$i\omega(x_AT(f)y_D)=i\omega([x_A,T(f)]y_D)$'' uses $T(f)\Omega=0$ twice
($\omega(x_AT(f)y_D)=\langle x_A^*\Omega,T(f)y_D\Omega\rangle$ and $T(f)^*\Omega=0$), and the conclusion
$g\rightsquigarrow-f(1-\chi)$ is right.  In the $\xi$-chart the two representatives are
$\widetilde g=-4\sinh^2\frac\xi2\mathbf 1_{\xi>0}$ and $\widetilde g_2=+4\sinh^2\frac\xi2\mathbf 1_{\xi<0}$, which
differ by $\widetilde g_2-\widetilde g=4\sinh^2\frac\xi2$, the $\xi$-image of the M\"obius field $x^2/L$
(indeed $\frac{x^2}{L(L-x)}=4\sinh^2\frac\xi2$).  Both have the same Fourier transform,
$\widehat{\widetilde g_2}(k)=-\widehat{\widetilde g}(-k)=\frac{-2i}{k(k^2+1)}=\widehat{\widetilde g}(k)$,
hence the same spectral density --- a consistency check the note does not make but which holds.
\end{verdictok}
\begin{lemma}[N7 eq.~(6),(7)]\label{lem:gt}
$\widetilde g(\xi)=-4\sinh^2\frac\xi2$ for $\xi>0$, and
$\widehat{\widetilde g}(k)=\lim_{\eps\downarrow0}\int_0^\infty e^{-(ik+\eps)\xi}\widetilde g\,d\xi\bigr|_{\text{cont.\ from }\eps>1}
=\frac{-2i}{k(k^2+1)}$, $|\widehat{\widetilde g}(k)|^2=\frac4{k^2(k^2+1)^2}$.
\end{lemma}
\lstep{1}{1}{$\widetilde g=g/\beta=\frac{-x^2/L}{L-x}=-\frac{L^2(1-e^{-\xi})^2}{L\cdot Le^{-\xi}}=-(e^{\xi/2}-e^{-\xi/2})^2=-4\sinh^2\frac\xi2$.}
\lby{Lemma~\ref{lem:W}$\langle1\rangle3$, eq.~(2), $f(x)=-x^2/L$}
\lstep{1}{2}{$-4\sinh^2\frac\xi2=-e^\xi-e^{-\xi}+2$.}
\lby{$4\sinh^2\frac\xi2=2\cosh\xi-2$}
\lstep{1}{3}{For $\eps>1$, $\int_0^\infty e^{-(ik+\eps)\xi}e^{a\xi}d\xi=(ik+\eps-a)^{-1}$ converges absolutely for
$a\in\{1,-1,0\}$; hence the integral equals $-\frac1{ik+\eps-1}-\frac1{ik+\eps+1}+\frac2{ik+\eps}$, a rational
function of $\eps$ regular at $\eps=0$.}
\lby{$\langle1\rangle2$, elementary integration, $\operatorname{Re}(ik+\eps-a)=\eps-a>0$}
\lstep{1}{4}{At $\eps=0$: $-2\bigl[\frac12\bigl(\frac1{ik-1}+\frac1{ik+1}\bigr)-\frac1{ik}\bigr]
=-2\bigl[\frac{-ik}{k^2+1}+\frac ik\bigr]=-2\cdot\frac{i}{k(k^2+1)}$.}
\lby{$\langle1\rangle3$, $\frac12\bigl(\frac1{ik-1}+\frac1{ik+1}\bigr)=\frac{ik}{(ik)^2-1}=\frac{-ik}{k^2+1}$ and
$\frac1{ik}=-\frac ik$; then $\frac1k-\frac k{k^2+1}=\frac1{k(k^2+1)}$}
\lqed{1}{5}\lby{$\langle1\rangle4$; $|{-2i}/(k(k^2+1))|^2=4/(k^2(k^2+1)^2)$}
\begin{verdictok}
Eqs.~(6) and (7) verified exactly as printed, including the analytic-continuation prescription, which is
unambiguous because the $\eps$-integral is a \emph{rational} function of $\eps$ (step $\langle1\rangle3$): the
continuation from $\eps>1$ to $\eps=0$ is unique.  This is worth saying in the note --- ``analytic continuation''
alone would not fix a function up to an entire additive term; rationality does.
\end{verdictok}

\begin{lemma}[N7 eq.~(8),(9)]\label{lem:D1}
With $V_s\phi(z)=\phi(h_s^{-1}(z))/\sqrt{h_s'(h_s^{-1}(z))}$, $h_s(x)=\frac{Lx}{L+sx}$, $f=\frac d{ds}h_s|_0=-x^2/L$:
$\mathfrak g:=-\frac{dW_s}{ds}|_0=(g\partial_x+\frac12g')E_D$ and $D^{(1)}=\mathfrak g^*Q+Q\mathfrak g$; and in the
$\xi$-frame $\widetilde D^{(1)}(\xi,\xi')=-\frac i{2\pi}\frac{\sinh\frac\xi2\sinh\frac{\xi'}2}{\sinh^2\frac{\xi-\xi'}2}$
for $\xi<0<\xi'$.
\end{lemma}
\lstep{1}{1}{$h_s(x)=x+sf(x)+O(s^2)$ with $f=-x^2/L$, hence $h_s^{-1}(z)=z-sf(z)+O(s^2)$ and
$\sqrt{h_s'(h_s^{-1}(z))}=1+\frac s2f'(z)+O(s^2)$.}
\lby{$\frac{d}{ds}\frac{Lx}{L+sx}\big|_0=-\frac{Lx^2}{L^2}=-\frac{x^2}L$; inverse-function expansion; $h_s'=1+sf'+O(s^2)$}
\lstep{1}{2}{$(V_s\phi)(z)=\phi(z)-s\bigl(f\phi'+\frac12f'\phi\bigr)(z)+O(s^2)$, so
$-\frac{dV_s}{ds}\big|_0=f\partial+\frac12f'$; restricted to $D$ this is $\mathfrak g$ of eq.~(8) with $g=f|_D$.}
\lby{$\langle1\rangle1$ and $W_s=1_{H_A}\oplus V_s$}
\lstep{1}{3}{$D^{(1)}=\frac d{ds}(Q-W_s^*QW_s)|_0=-\bigl(\tfrac{dW^*}{ds}Q+Q\tfrac{dW}{ds}\bigr)\big|_0
=\mathfrak g^*Q+Q\mathfrak g$; it is self-adjoint.}
\lby{$\langle1\rangle2$, product rule, $Q=Q^*$}
\lstep{1}{4}{$D^{(1)}$ vanishes on $A\times A$ and on $D\times D$.}
\lby{$h_s$ is a global M\"obius map and $P_+$ is M\"obius invariant; $h_s(D)\subset D$, so for $\phi,\phi'$
supported in $D$, $\langle\phi,W_s^*P_+W_s\phi'\rangle=\langle V_s\phi,P_+V_s\phi'\rangle=\langle\phi,P_+\phi'\rangle$,
and $W_s=1$ on $A$}
\lstep{1}{5}{Setting $u=-x=L(e^{-\xi}-1)$, $v=y=L(1-e^{-\xi'})$: $u=-2Le^{-\xi/2}\sinh\frac\xi2$,
$v=2Le^{-\xi'/2}\sinh\frac{\xi'}2$, $u+v=L(e^{-\xi}-e^{-\xi'})=-2Le^{-(\xi+\xi')/2}\sinh\frac{\xi-\xi'}2$.}
\lby{$e^{-\xi}-1=-2e^{-\xi/2}\sinh\frac\xi2$, $1-e^{-\xi'}=2e^{-\xi'/2}\sinh\frac{\xi'}2$ and Lemma~\ref{lem:W}$\langle1\rangle2$}
\lqed{1}{6}
\lby{$\widetilde D^{(1)}=\sqrt{\beta\beta'}\,\frac i{2\pi L}\frac{uv}{(u+v)^2}$ with $\sqrt{\beta\beta'}=Le^{-(\xi+\xi')/2}$
and $\langle1\rangle5$: $\frac{uv}{(u+v)^2}=\frac{-4L^2e^{-(\xi+\xi')/2}\sinh\frac\xi2\sinh\frac{\xi'}2}{4L^2e^{-(\xi+\xi')}\sinh^2\frac{\xi-\xi'}2}
=-e^{(\xi+\xi')/2}\frac{\sinh\frac\xi2\sinh\frac{\xi'}2}{\sinh^2\frac{\xi-\xi'}2}$}
\begin{verdictok}
Eqs.~(8) and (9) verified, \emph{including the sign of $\mathfrak g$}, which is the one place where a slip would
flip the sign of $D^{(1)}$ (and would be invisible in $|D^{(1)}|^2$, so it does not endanger $f_2$ anyway).
The claimed decay $e^{-|\xi|/2}e^{-\xi'/2}$ and boundedness are correct: for $\xi<0<\xi'$,
$|\sinh\frac\xi2\sinh\frac{\xi'}2|\le\sinh^2\frac{\xi'-\xi}4\le\sinh^2\frac{\xi'-\xi}2$ by AM--GM, so
$|\widetilde D^{(1)}|\le\frac1{2\pi}$.  Step $\langle1\rangle4$ (``on $A\times A$ and $D\times D$ both sides
vanish'') is asserted without proof in the note; the M\"obius-invariance argument above supplies it.
\end{verdictok}
\section{Section 3 of Note 7: the modular quadratic forms}\label{sec:lem31}

\begin{lemma}[N7 \S3.1 and Lemma 3.1: all identities]\label{lem:std}
In the finite-dimensional standard form of N7 \S3.1: (a) $\Delta=\rho\otimes\bar\rho^{-1}$, $S=J\Delta^{1/2}$,
$S^*=\Delta^{1/2}J$; (b) $\eta_{i\bar\jmath}=d\rho_{ij}/\sqrt{p_j}$ is the unique vector with
$\Tr(d\rho X)=\langle\eta,X\Omega\rangle$; (c) eq.~(11) holds, i.e.\
$g_B=2\langle\eta,(\Delta+1)^{-1}\eta\rangle$ and $g_{\mathrm{KM}}=\langle\eta,\frac{\log\Delta}{\Delta-1}\eta\rangle$;
(d) $\eta=i(S^*-1)A\Omega$ and consequences (1)--(3); (e) $d\mu(\lambda^{-1})=\lambda\,d\mu(\lambda)$;
(f) $w(\lambda^{-1})=w(\lambda)/\lambda$ for both weights.
\end{lemma}
\lstep{1}{1}{(a) $J(|i\rangle\otimes|\bar\jmath\rangle)=|j\rangle\otimes|\bar\imath\rangle$ antilinearly and
$\Delta^{1/2}|i\bar\jmath\rangle=\sqrt{p_i/p_j}|i\bar\jmath\rangle$ give
$J\Delta^{1/2}(X\Omega)=X^*\Omega$ for $X=|a\rangle\langle b|$: $X\Omega=\sqrt{p_b}|a\bar b\rangle\mapsto
\sqrt{p_b}\sqrt{p_a/p_b}|b\bar a\rangle=\sqrt{p_a}|b\bar a\rangle=X^*\Omega$.  $S^*=\Delta^{1/2}J$ because
$J^*=J=J^{-1}$ and $J\Delta J=\Delta^{-1}$.}
\lby{direct computation; $\{|a\rangle\langle b|\}$ spans $\cM$}
\lstep{1}{2}{(b) $\Tr(d\rho X)=\sum_{ij}(d\rho)_{ji}X_{ij}$ and $\langle\eta,X\Omega\rangle=\sum_{ij}\overline{\eta_{i\bar\jmath}}X_{ij}\sqrt{p_j}$;
equating and using $d\rho=d\rho^*$ gives $\eta_{i\bar\jmath}=d\rho_{ij}/\sqrt{p_j}$.  Uniqueness: $\Omega$ is
cyclic for $\cM$.}
\lby{comparison of coefficients; $\cM\Omega$ dense}
\lstep{1}{3}{(c) $\langle\eta,(\Delta+1)^{-1}\eta\rangle=\sum_{ij}\frac{|d\rho_{ij}|^2}{p_j}\frac{1}{p_i/p_j+1}
=\sum_{ij}\frac{|d\rho_{ij}|^2}{p_i+p_j}$, which is $\frac12g_B$; and
$\langle\eta,\frac{\log\Delta}{\Delta-1}\eta\rangle=\sum\frac{|d\rho_{ij}|^2}{p_j}\frac{\log\lambda}{\lambda-1}
=\sum|d\rho_{ij}|^2\ell(p_i,p_j)=g_{\mathrm{KM}}$, $\lambda=p_i/p_j$.}
\lby{$\langle1\rangle2$, eq.~(10), and $p_j\ell(p_i,p_j)=\frac{\log\lambda}{\lambda-1}$ (the note's own identity,
correct)}
\lstep{1}{4}{The auxiliary identity of the note,
$(1\otimes\bar\rho^{1/2})(\rho\otimes1+1\otimes\bar\rho)^{-1}(1\otimes\bar\rho^{1/2})=(\Delta+1)^{-1}$, holds:
on $|i\bar\jmath\rangle$ both sides are $p_j/(p_i+p_j)$.}
\lby{spectral computation}
\lstep{1}{5}{(d) $\varphi(X)=i\langle A\Omega,X\Omega\rangle-i\langle X^*\Omega,A\Omega\rangle$ and
$\langle X^*\Omega,A\Omega\rangle=\langle J\Delta^{1/2}X\Omega,A\Omega\rangle=\langle JA\Omega,\Delta^{1/2}X\Omega\rangle
=\langle\Delta^{1/2}JA\Omega,X\Omega\rangle=\langle S^*A\Omega,X\Omega\rangle$, so
$\varphi(X)=\langle i(S^*-1)A\Omega,X\Omega\rangle$.}
\lby{$\langle Jx,y\rangle=\overline{\langle x,Jy\rangle}=\langle Jy,x\rangle$ for antiunitary $J$, and
$\Delta^{1/2}=(\Delta^{1/2})^*$; antilinearity of $\eta\mapsto\langle\eta,\cdot\rangle$ turns
$i\langle A\Omega,\cdot\rangle$ into $\langle-iA\Omega,\cdot\rangle$}
\lstep{1}{6}{(1): for $B=B^*\in\cM'$, $S^*$ is the Tomita operator of $\cM'$, so $S^*B\Omega=B^*\Omega=B\Omega$
and $(S^*-1)B\Omega=0$.  (2): $JX\Omega=\Delta^{1/2}X^*\Omega$ (apply $J$ to $J\Delta^{1/2}X^*\Omega=X\Omega$),
hence $S^*X\Omega=\Delta X^*\Omega$, and $A=A^*\in\cM$ gives $\eta=i(\Delta-1)A\Omega$; substituting in
$\langle1\rangle3$ gives eq.~(13).  (3): $(1+\Delta)^{-1/2}\Delta^{1/2}=(1+\Delta^{-1})^{-1/2}$ and
$J(1+\Delta)^{-1/2}J=(1+\Delta^{-1})^{-1/2}$ give $(1+\Delta)^{-1/2}\Delta^{1/2}J=J(1+\Delta)^{-1/2}$.}
\lby{Tomita--Takesaki for $\cM'$; $J\Delta J=\Delta^{-1}$; functional calculus}
\lstep{1}{7}{(e) $\mu=\sum_{ij}|A_{ij}|^2p_j\,\delta_{p_i/p_j}$; exchanging $i\leftrightarrow j$ and using
$|A_{ji}|=|A_{ij}|$ gives $\int F(\lambda)d\mu=\int\lambda F(1/\lambda)d\mu$, i.e.\ $d\mu(1/\lambda)=\lambda d\mu(\lambda)$.}
\lby{$A=A^*$}
\lstep{1}{8}{(f) $w_B(1/\lambda)=\frac{(1/\lambda-1)^2}{1/\lambda+1}=\frac{(1-\lambda)^2}{\lambda(1+\lambda)}=w_B(\lambda)/\lambda$;
$w_{\mathrm{KM}}(1/\lambda)=(1/\lambda-1)\log(1/\lambda)=\frac{(\lambda-1)\log\lambda}{\lambda}=w_{\mathrm{KM}}(\lambda)/\lambda$.
Combined with (e) this gives ``twice the contribution of $\{\lambda<1\}$''.}
\lby{algebra; $\int_{\lambda>1}w\,d\mu=\int_{\lambda<1}\lambda w(1/\lambda)d\mu=\int_{\lambda<1}w\,d\mu$ by (e),(f)}
\lqed{1}{9}\lby{$\langle1\rangle1$--$\langle1\rangle8$}
\begin{verdictok}
Every identity of N7 \S3.1 and Lemma~3.1, including all three consequences, the eigenbasis formula, the KMS
property and both weight symmetries, is verified.  The finite-dimensional inputs eq.~(10) (SLD quantum Fisher
information $=2\sum|d\rho_{ij}|^2/(p_i+p_j)$ with $-\log\Fid=\frac{s^2}8g_B+o(s^2)$, and the Kubo--Mori Hessian
of the relative entropy) are Note~5's Prop.~2.3, independently verified by [V3] and [R3].
\end{verdictok}
\subsection{A false boundedness claim}\label{sec:bound}
\begin{issue}[N7, sentence after eq.~(13)]\label{iss:bound}
``\emph{Both weights in \textup{(13)} satisfy $w(\lambda^{-1})=w(\lambda)/\lambda$, so each integral is twice the
contribution of $\{\lambda<1\}$; and both are bounded by $1+\lambda$, which is why $g_B\le4\|A\Omega\|^2$
(equality for pure states) and why the second variations are finite whenever $A\Omega$ is a vector.}''
\end{issue}
\lstep{1}{1}{$w_B(\lambda)=\frac{(\lambda-1)^2}{\lambda+1}\le1+\lambda$ for $\lambda\ge0$, with equality only at
$\lambda=0$; hence $g_B=2\int w_Bd\mu\le2\int(1+\lambda)d\mu=2(\|A\Omega\|^2+\|\Delta^{1/2}A\Omega\|^2)=4\|A\Omega\|^2$.}
\lby{$(\lambda-1)^2\le(1+\lambda)^2$ for $\lambda\ge0$; $\|\Delta^{1/2}A\Omega\|=\|J\Delta^{1/2}A\Omega\|=\|A^*\Omega\|=\|A\Omega\|$ for $A=A^*\in\cM$}
\lstep{1}{2}{$w_{\mathrm{KM}}(\lambda)=(\lambda-1)\log\lambda$ is \emph{not} bounded by $1+\lambda$:
$w_{\mathrm{KM}}(\lambda)/(1+\lambda)\to\infty$ as $\lambda\to\infty$.  Explicitly at $\lambda=e^2=7.389$:
$w_{\mathrm{KM}}=12.7780>8.3891=1+\lambda$; at $\lambda=100$: $455.91>101$.}
\lby{numerical evaluation, \texttt{numerics/v5chk.py}}
\lstep{1}{3}{Consequently $g_{\mathrm{KM}}=\int(\lambda-1)\log\lambda\,d\mu$ can diverge for $A\Omega\in\mathcal H$:
finiteness needs the extra moment $\int\lambda\log\lambda\,d\mu<\infty$, i.e.\ $A\Omega$ in the form domain of
$\Delta\log\Delta$, which is exactly the hypothesis Lemma~3.2 has to impose separately.}
\lby{$\langle1\rangle2$; e.g.\ $d\mu$ with density $\propto\lambda^{-2}(\log\lambda)^{-2}$ on $\lambda>e$ has
$\int d\mu<\infty$, $\int\lambda\log\lambda\,d\mu=\infty$}
\lqed{1}{4}\lby{$\langle1\rangle1$--$\langle1\rangle3$}
\begin{verdictbreak}
The clause ``\emph{and both are bounded by $1+\lambda$ \dots\ and why the second variations are finite whenever
$A\Omega$ is a vector}'' is \textbf{false} for the Kubo--Mori weight, and the finiteness conclusion it draws is
false for $g_{\mathrm{KM}}$.  It is a side remark and does not enter any computation --- $g_{\mathrm{KM}}$ is
evaluated by an explicit convergent integral in \S4.3 and \S5 --- so no result of the note falls.  But it is
inconsistent with Lemma~3.2, which (correctly) imposes ``$\eta$ in the form domain of $\log\Delta/(\Delta-1)$''
precisely because boundedness fails.  Repair: replace by ``the Bures weight is bounded by $1+\lambda$, whence
$g_B\le4\|A\Omega\|^2$; the Kubo--Mori weight is not, and $g_{\mathrm{KM}}<\infty$ requires the extra condition
$A\Omega\in\mathcal Q(\Delta\log\Delta)$.''  Also ``equality for pure states'' is loose: a pure state is not
faithful and has no standard form of the type used.
\end{verdictbreak}

\subsection{$\|T(g)\Omega\|=\infty$, and what it costs}\label{sec:normdiv}
\begin{proposition}\label{prop:infnorm}
For $\widetilde g=-4\sinh^2\frac\xi2\mathbf 1_{\xi>0}$ the formal spectral integral
$\int d\mu=\frac1{2\pi}\int_{\mathbb R}\widehat W(k)|\widehat{\widetilde g}(k)|^2dk$ diverges, and no representative
$g$ of the tangent functional with $g=f$ on a neighbourhood of $\overline D$ lies in $\cA(I)$.
\end{proposition}
\lstep{1}{1}{$\widehat W(0)=\frac{c}{48\pi^2}>0$ and $|\widehat{\widetilde g}(k)|^2=\frac4{k^2(k^2+1)^2}\sim\frac4{k^2}$
as $k\to0$, so the integrand behaves like $\frac{c}{24\pi^3}k^{-2}$, which is not integrable at $0$ (the
singularity is even and positive: no principal value helps).}
\lby{Lemma~\ref{lem:What7} below and Lemma~\ref{lem:gt}; numerically $\widehat W(0)=0.00211085799255=1/(48\pi^2)$
at $c=1$}
\lstep{1}{2}{Adding to $g$ any of $1,x,x^2$ does not change $T(g)\Omega$ (they annihilate $\Omega$), so the
divergence cannot be removed inside the $\xi$-chart.}
\lby{N7 \S1, M\"obius invariance of $\Omega$}
\lstep{1}{3}{The $\xi$-chart covers only $x<L$; a representative with $T(g)\Omega\in\mathcal H$ must be obtained by
cutting $g$ off \emph{beyond} $L$, i.e.\ outside the chart, and then $\operatorname{supp}g\not\subset I$, so
$T(g)$ is not affiliated with $\cA(I)$.}
\lby{$\xi=-\log\frac{L-x}L$ maps $(-\infty,L)$ onto $\mathbb R$; Lemma~2.1 forces $g=f$ on a neighbourhood of
$\overline D=[0,L]$, hence $g\ne0$ on $(L,L+\delta)$}
\lqed{1}{4}\lby{$\langle1\rangle1$--$\langle1\rangle3$}
\begin{verdictgap}
Confirms [V6].  Two consequences the note should draw explicitly.
(i) \textbf{Lemma~3.1(2) does not apply to $A=T(g)$}: it assumes $A\in\cM$, and step $\langle1\rangle3$ shows no
admissible representative is.  Since \S4.2--4.3 use precisely eq.~(13) --- the $A\in\cM$ formula --- the general
route rests on a formula whose hypothesis fails.  Lemma~3.1(3) is the correct general statement, and it is
representative-independent by Lemma~3.1(1) (two admissible $g$'s differ by a function supported in $(L,\infty)$,
so $T(g_1)-T(g_2)$ is self-adjoint and affiliated with $\cA((L,\infty))=\cA(I)'$ by Haag duality).
(ii) \textbf{The two formulas do agree here}, which is why the answer survives: for $A\in\cM$ one has
$g_B=4\|\chi\|^2-4\operatorname{Re}\langle J\chi,\chi\rangle=4\int\frac{1-\lambda}{1+\lambda}d\mu$, and the KMS
symmetry (Lemma~\ref{lem:std}(e)) turns this into $2\int\frac{(\lambda-1)^2}{1+\lambda}d\mu$; and the note's
$d\mu$, although it is not the spectral measure of a vector, does satisfy the KMS relation, because
$\widehat W(-k)=e^{2\pi k}\widehat W(k)$ and $|\widehat{\widetilde g}(-k)|=|\widehat{\widetilde g}(k)|$.  Note that
$4\int\frac{1-\lambda}{1+\lambda}d\mu$ is only conditionally convergent at $k=0$ (odd integrand $\sim -1/k$) while
the symmetrized $2\int\frac{(\lambda-1)^2}{1+\lambda}d\mu$ is absolutely convergent: the symmetrization is what
makes the formal computation finite.  This is the honest description of \S4, and it is a \emph{gap}, not a break.
\end{verdictgap}
\subsection{Lemma 3.2: is it well posed?}
\begin{verdictgap}
The statement is \emph{almost} well posed; three defects.
\begin{enumerate}[leftmargin=1.6em]
\item ``\emph{a self-adjoint $A$ affiliated with a von Neumann algebra containing $\cM$}'' is a vacuous
hypothesis: every self-adjoint operator on $\mathcal H$ is affiliated with $B(\mathcal H)\supseteq\cM$.  What is
presumably meant, and what the proof will need, is $A$ affiliated with $\cA(\widehat I)$ for some interval
$\widehat I\supset I$, so that locality and the Reeh--Schlieder property are available.  As written the
hypothesis carries no information, and in particular gives no control of $S^*A\Omega$.
\item ``\emph{$o(s)$ is controlled in the norm of the tangent space}'' is not a definition.  The tangent space
here is $\cM_*$ (or its completion in the relevant metric); which norm, and uniformity over which set of $X$,
must be specified --- this is exactly the point where the finite-dimensional Taylor expansion is transferred.
\item The condition ``$\eta$ in the form domain of $\log\Delta/(\Delta-1)$'' is the right one for
$g_{\mathrm{KM}}$; for $g_B$ nothing is needed beyond $\eta\in\mathcal H$, since $(1+\Delta)^{-1}$ is bounded.
Worth splitting.
\end{enumerate}
\textbf{Consistency with Lemma 3.1}: good.  Lemma~3.2 refers to eq.~(11) (the $\eta$-form), not eq.~(13) (the
$A\in\cM$ spectral form), so it is compatible with $A\notin\cM$; and it requires $\Omega\in D(A)$, which is
exactly what Prop.~\ref{prop:infnorm} says fails for the zero-collar $\xi$-frame representative.  So Lemma~3.2,
even once proved, will \emph{not} directly license \S4.2--4.3; it licenses Lemma~3.1(3) for a compactly
supported $g$, and the passage from there to the $\xi$-frame integrals is a separate (missing) step.  The note's
own caveat in \S3.3 now says this; the caveat should be repeated at eq.~(15), where the damage is done.
The corrected proof route (sup formula, Alberti--Uhlmann 2002 Thm.~2(3)) is right --- the previously stated
$\Fid=\langle\Omega,\Delta_{\omega_s,\omega}^{1/2}\Omega\rangle$ is indeed $Q_{1/2}$, not the fidelity [R3];
I confirm that refutation and add that the note's remaining phrase ``\emph{and the finite-dimensional variational
characterization of the Bures metric}'' is doing a lot of unstated work.
\end{verdictgap}

\section{Section 4 of Note 7: the modular-frame evaluation}\label{sec:step3}

\begin{lemma}[N7 Lemma 4.1]\label{lem:What7}
$\displaystyle\Wh(k)=\int_{\mathbb R}e^{-iku}W(u)du=\frac{c}{24\pi}\frac{k(k^2+1)}{e^{2\pi k}-1}$, with
$\Wh(-k)=e^{2\pi k}\Wh(k)$, $\Wh\ge0$, $\Wh(0)=\frac c{48\pi^2}$, and $\Wh(k)\simeq\frac c{24\pi}|k|^3\theta(-k)$
as $|k|\to\infty$.
\end{lemma}
\lstep{1}{1}{$W$ is the boundary value from $\operatorname{Im}z<0$ of $\frac{c}{128\pi^2}\sinh^{-4}(z/2)$, whose
poles are the fourth-order poles $z\in2\pi i\mathbb Z$; the strip $-2\pi<\operatorname{Im}z<0$ is pole free, so the
contour may be moved to $\operatorname{Im}u=-\pi$.}
\lby{Lemma~\ref{lem:W}; $\sinh(z/2)=0\Leftrightarrow z/2\in i\pi\mathbb Z$}
\lstep{1}{2}{The two vertical segments $\operatorname{Re}u=\pm R$, $-\pi\le\operatorname{Im}u\le0$ contribute $0$ in
the limit $R\to\infty$.}
\lby{$|\sinh\frac{\pm R+iy}2|^2=\sinh^2\frac R2\cos^2\frac y2+\cosh^2\frac R2\sin^2\frac y2\ge\sinh^2\frac R2$, so
$|\sinh^{-4}|\le\sinh^{-4}\frac R2=O(e^{-2R})$, while $|e^{-iku}|\le\max(1,e^{\pi k})$ and the segment length is $\pi$}
\lstep{1}{3}{On $u=v-i\pi$: $\sinh\frac{v-i\pi}2=\sinh\frac v2\cos\frac\pi2-i\cosh\frac v2\sin\frac\pi2=-i\cosh\frac v2$,
so $\sinh^4=\cosh^4\frac v2$, and $e^{-ik(v-i\pi)}=e^{-ikv}e^{-\pi k}$.  Hence
$\Wh(k)=\frac{c}{128\pi^2}e^{-\pi k}\int e^{-ikv}\cosh^{-4}\frac v2\,dv$.}
\lby{$\langle1\rangle1$, $\langle1\rangle2$, $\sinh(a-i\pi/2)=-i\cosh a$}
\lstep{1}{4}{$\int_{\mathbb R}e^{-ikv}\cosh^{-2h}\frac v2\,dv=\frac{2^{2h}}{\Gamma(2h)}|\Gamma(h+ik)|^2$; for $h=2$
this is $\frac{16}6|\Gamma(2+ik)|^2=\frac83|\Gamma(2+ik)|^2$.}
\lby{standard beta-integral (Gradshteyn--Ryzhik 3.985.1); verified numerically below}
\lstep{1}{5}{$|\Gamma(2+ik)|^2=|1+ik|^2|\Gamma(1+ik)|^2=(1+k^2)\frac{\pi k}{\sinh\pi k}$.}
\lby{$\Gamma(2+ik)=(1+ik)\Gamma(1+ik)$ and the reflection formula $|\Gamma(1+ik)|^2=\pi k/\sinh\pi k$}
\lstep{1}{6}{$\Wh(k)=\frac{c}{128\pi^2}e^{-\pi k}\cdot\frac83\cdot\frac{\pi k(k^2+1)}{\sinh\pi k}
=\frac{ck(k^2+1)}{48\pi}\cdot\frac{e^{-\pi k}}{\sinh\pi k}=\frac{c}{24\pi}\frac{k(k^2+1)}{e^{2\pi k}-1}$.}
\lby{$\langle1\rangle3$--$\langle1\rangle5$, $\frac{e^{-\pi k}}{\sinh\pi k}=\frac2{e^{2\pi k}-1}$}
\lstep{1}{7}{$\Wh\ge0$ (numerator and denominator have the same sign); $\Wh(0)=\frac c{24\pi}\cdot\frac1{2\pi}=\frac c{48\pi^2}$;
$\Wh(-k)/\Wh(k)=\frac{-k}{e^{-2\pi k}-1}\cdot\frac{e^{2\pi k}-1}k=e^{2\pi k}$.}
\lby{$\langle1\rangle6$, l'H\^opital at $k=0$}
\lstep{1}{8}{As $k\to-\infty$, $e^{2\pi k}-1\to-1$ and $\Wh\to\frac{c}{24\pi}|k|^3$; as $k\to+\infty$, $\Wh\to0$
exponentially.  The flat-space check: $\int e^{-ipu}(u-i0)^{-4}du=2\pi i\,\theta(-p)\frac{(-ip)^3}{3!}=\frac{2\pi}{3!}|p|^3\theta(-p)$,
so $\frac c{8\pi^2}\cdot\frac{2\pi}{6}|p|^3=\frac{c}{24\pi}|p|^3$ for $p<0$, matching.}
\lby{$\langle1\rangle6$; $(u-i0)^{-n}=\frac{(-1)^{n-1}}{(n-1)!}\partial_u^{n-1}(u-i0)^{-1}$,
$\int e^{-ipu}(u-i0)^{-1}du=2\pi i\theta(-p)$ (close in the upper half plane), and
$\int e^{-ipu}f^{(m)}=(ip)^m\int e^{-ipu}f$}
\lqed{1}{9}\lby{$\langle1\rangle6$--$\langle1\rangle8$}
\begin{verdictok}
Lemma~4.1 is correct in every detail: the contour shift, the identity $\sinh\frac{v-i\pi}2=-i\cosh\frac v2$, the
$\cosh^{-2h}$ transform, $|\Gamma(2+ik)|^2$, the KMS relation and the flat-space cross-check.  Numerically the
contour-shifted quadrature reproduces $\frac{c}{24\pi}\frac{k(k^2+1)}{e^{2\pi k}-1}$ to 14 digits at
$k=-1.3,0.25,1.0$; $|\Gamma(2+ik)|^2=\pi k(1+k^2)/\sinh\pi k$ to 12 digits; and the flat-space transform of
$(u-i0)^{-4}$ agrees with $\frac{2\pi}6|p|^3\theta(-p)$ (5.14486 vs 5.14488 at $p=-1.7$; $-1.6\times10^{-5}$ vs
$0$ at $p=+0.9$).  \textbf{The flat-space check is an independent confirmation of eq.~(1)}: the coefficient
$\frac c{24\pi}$ of $|k|^3$ is the same one Theorem~\ref{thm:TT} produces.
\end{verdictok}
\subsection{The KMS sign in \S4.2}\label{sec:kmssign}
\begin{issue}[N7 eq.~(15) and the sentence following it]
``\emph{the sign being fixed by the KMS asymmetry $d\mu(-u)=e^ud\mu(u)$ of \textup{(13)}, which is
$\Wh(-k)=e^{2\pi k}\Wh(k)$}.''
\end{issue}
\lstep{1}{1}{The KMS relation $d\mu(-u)=e^ud\mu(u)$ of eq.~(13) is a property of the spectral measure of $A\Omega$
for $A\in\cM$ (Lemma~\ref{lem:std}(e)); the proof uses $|A_{ij}|^2p_i=\lambda|A_{ij}|^2p_j$, i.e.\ that $A$ is a
matrix over the \emph{same} algebra whose density matrix defines $\Delta$.}
\lby{Lemma~\ref{lem:std}(e) and its proof}
\lstep{1}{2}{$A=T(g)\notin\cM$ and $A\Omega\notin\mathcal H$.}
\lby{Prop.~\ref{prop:infnorm}}
\lstep{1}{3}{Hence the argument of the issue is circular in the following precise sense: it invokes a property
that holds for vectors $A\Omega$, $A\in\cM$, in order to determine the modular flow direction for an $A$ that is
not in $\cM$ and whose $A\Omega$ is not a vector.}
\lby{$\langle1\rangle1$, $\langle1\rangle2$}
\lstep{1}{4}{The sign is nevertheless the one the note uses: by Lemma~\ref{lem:bw},
$\Delta^{it}\widetilde T_0(\xi)\Delta^{-it}=\widetilde T_0(\xi-2\pi t)$, so
$\langle T(g)\Omega,\Delta^{it}T(g)\Omega\rangle=\iint\gt\gt'W(\xi-\xi'+2\pi t)
=\frac1{2\pi}\int\Wh(k)|\gh(k)|^2e^{2\pi ikt}dk$, and comparing with $\int e^{iut}d\mu(u)$ gives $u=2\pi k$.}
\lby{Lemma~\ref{lem:bw}; $\Delta^{\pm it}\Omega=\Omega$; $W(u)=\frac1{2\pi}\int\Wh(k)e^{iku}dk$;
$\int\gt(\xi)e^{ik\xi}d\xi=\overline{\gh(k)}$ since $\gt$ is real}
\lqed{1}{5}\lby{$\langle1\rangle3$, $\langle1\rangle4$}
\begin{verdictgap}
The \emph{result} of \S4.2 is correct ($u=\log\lambda=2\pi k$), but the justification offered is not valid: it
uses a KMS symmetry that is a theorem only for $A\in\cM$, in a situation where $A\notin\cM$.  Replace by
Lemma~\ref{lem:bw} (Bisognano--Wichmann plus Haag duality), which settles the sign with no reference to the
tangent vector at all.  The note's own added caveat --- ``\emph{the total mass $\int d\mu$ diverges \dots\ and
$T(g)\Omega$ is not a vector for this representative}'' --- is correct and confirmed by
Prop.~\ref{prop:infnorm}$\langle1\rangle1$; it should be sharpened to say that eq.~(13), not merely the total
mass, is what is unavailable.
\end{verdictgap}

\subsection{The two integrals of \S4.3}
\begin{lemma}[N7 eqs.~(16)--(19)]\label{lem:ints}
With $d\mu(u)=\frac1{2\pi}\Wh(k)|\gh(k)|^2dk$, $u=2\pi k$, $\lambda=e^{2\pi k}$:
$g_B=\frac c{6\pi^2}\int_{\mathbb R}\frac{\tanh\pi k}{k(k^2+1)}dk=\frac{2c}{3\pi^2}$ and
$g_{\mathrm{KM}}=\frac c{6\pi}\int_{\mathbb R}\frac{dk}{k^2+1}=\frac c6$; and
$\int_0^\infty\frac{\tanh\pi k}{k(k^2+1)}dk=2$.
\end{lemma}
\lstep{1}{1}{$g_B=2\int\frac{(\lambda-1)^2}{\lambda+1}d\mu=\frac1\pi\int dk\,\Wh|\gh|^2\frac{(e^{2\pi k}-1)^2}{e^{2\pi k}+1}$.}
\lby{eq.~(13) and $d\mu=\frac1{2\pi}\Wh|\gh|^2dk$; $2\cdot\frac1{2\pi}=\frac1\pi$}
\lstep{1}{2}{$\frac1\pi\cdot\frac c{24\pi}\frac{k(k^2+1)}{e^{2\pi k}-1}\cdot\frac{(e^{2\pi k}-1)^2}{e^{2\pi k}+1}\cdot\frac4{k^2(k^2+1)^2}
=\frac{c}{6\pi^2}\frac{\tanh\pi k}{k(k^2+1)}$.}
\lby{$\frac{4}{24\pi^2}=\frac1{6\pi^2}$; $\frac{k(k^2+1)}{k^2(k^2+1)^2}=\frac1{k(k^2+1)}$;
$\frac{e^{2\pi k}-1}{e^{2\pi k}+1}=\tanh\pi k$}
\lstep{1}{3}{$g_{\mathrm{KM}}=\int(\lambda-1)\log\lambda\,d\mu=\frac1{2\pi}\int dk\,\Wh|\gh|^2(e^{2\pi k}-1)2\pi k
=\frac c{6\pi}\int\frac{dk}{k^2+1}=\frac c6$.}
\lby{$\frac{2\pi}{2\pi}\cdot\frac{c}{24\pi}\cdot4\cdot\frac{k\cdot k(k^2+1)}{k^2(k^2+1)^2}=\frac{c}{6\pi}\frac1{k^2+1}$;
$\int_{\mathbb R}\frac{dk}{k^2+1}=\pi$}
\lstep{1}{4}{$\tanh\pi k=\sum_{n\ge0}\frac{2k/\pi}{k^2+(n+\frac12)^2}$.}
\lby{$\tanh z=\sum_{n\ge0}\frac{2z}{z^2+(n+\frac12)^2\pi^2}$ (Mittag-Leffler) at $z=\pi k$}
\lstep{1}{5}{$\int_0^\infty\frac{dk}{(k^2+a^2)(k^2+b^2)}=\frac\pi{2ab(a+b)}$ for $a,b>0$.}
\lby{partial fractions $\frac1{(k^2+a^2)(k^2+b^2)}=\frac1{b^2-a^2}\bigl(\frac1{k^2+a^2}-\frac1{k^2+b^2}\bigr)$ and
$\int_0^\infty\frac{dk}{k^2+a^2}=\frac\pi{2a}$; the case $a=b$ by continuity}
\lstep{1}{6}{$\int_0^\infty\frac{\tanh\pi k}{k(k^2+1)}dk=\sum_{n\ge0}\frac2\pi\cdot\frac\pi{2(n+\frac12)(n+\frac32)}
=\sum_{n\ge0}\frac4{(2n+1)(2n+3)}=2$.}
\lby{$\langle1\rangle4$, $\langle1\rangle5$ with $a=n+\frac12$, $b=1$; interchange of $\sum$ and $\int$ by Tonelli
(all terms $\ge0$); $\sum_{n\ge0}\frac4{(2n+1)(2n+3)}=2\sum_{n\ge0}\bigl(\frac1{2n+1}-\frac1{2n+3}\bigr)=2$}
\lqed{1}{7}\lby{$\langle1\rangle2$, $\langle1\rangle6$ (the integrand is even, so $\int_{\mathbb R}=4$),
$g_B=\frac c{6\pi^2}\cdot4=\frac{2c}{3\pi^2}$, and $\langle1\rangle3$}
\begin{verdictok}
Eqs.~(16)--(19) verified, every factor.  $f_2=g_B/8=\frac{c}{12\pi^2}=0.0084434\,c$ and
$s_2=g_{\mathrm{KM}}/2=\frac c{12}$, so eq.~(19) is right, as is $s_2/f_2=\pi^2$.  Numerically
$\int_0^\infty\frac{\tanh\pi k}{k(k^2+1)}dk=2.0$ to 15 digits.  The remark that ``\emph{the pole of $|\gh|^2$ at
$k=0$ \dots\ is cancelled by $\tanh\pi k$ and by $k$ respectively}'' is exactly right and is what makes the
divergence of Prop.~\ref{prop:infnorm} harmless; ``\emph{the kink of $\gt$ at the touching point gives the
$k^{-3}$ decay}'' is also right ($\gt\in C^{1,1}$ at $0$ $\Rightarrow$ $\gh=O(k^{-3})$).
\end{verdictok}
\subsection{\S4.4: is the universality argument valid?}
\begin{proposition}[Linearity in $c$]\label{prop:lin}
Fix the geometry $(A,D,I)$ and a compactly supported representative $g$.  Then, granting Lemma~3.2, $g_B$ and
$g_{\mathrm{KM}}$ depend on the net only through the connected two-point function of $T$, hence are exactly
proportional to $c$; and $J\,T(g)\,J=T(g\circ r)$ with $r(x)=2L-x$.
\end{proposition}
\lstep{1}{1}{$g_B=4\|\chi\|^2-4\operatorname{Re}\langle J\chi,\chi\rangle$ with $\chi=(1+\Delta)^{-1/2}T(g)\Omega$.}
\lby{Lemma~\ref{lem:std}(d)(3); $\|(1-J)\chi\|^2=2\|\chi\|^2-2\operatorname{Re}\langle J\chi,\chi\rangle$, using
$\|J\chi\|=\|\chi\|$ and $\langle\chi,J\chi\rangle=\overline{\langle J\chi,\chi\rangle}$}
\lstep{1}{2}{$\|\chi\|^2=\int\frac{d\mu(\lambda)}{1+\lambda}$ and $\langle J\chi,\chi\rangle=\int\frac{\lambda\,d\mu(\lambda)}{1+\lambda}$,
both determined by the spectral measure of $T(g)\Omega$ with respect to $\Delta$.}
\lby{functional calculus; for the second, $J(1+\Delta)^{-1/2}=(1+\Delta^{-1})^{-1/2}J$ and, in the eigenbasis,
$\langle J\chi,\chi\rangle=\sum|A_{ij}|^2\sqrt{p_ip_j}(1+\lambda^{-1})^{-1/2}(1+\lambda)^{-1/2}=\int\frac{\lambda}{1+\lambda}d\mu$}
\lstep{1}{3}{That spectral measure is the Fourier transform of $t\mapsto\langle T(g)\Omega,\Delta^{it}T(g)\Omega\rangle$,
which by Lemma~\ref{lem:bw} is $\iint g\,g'\,\langle T T\rangle_{\mathrm{conn}}$ evaluated on translated
arguments; it is therefore linear in the constant of eq.~(1), i.e.\ in $c$.  $\Delta$, $J$ and the geometric
action are independent of $c$ (BGL93).}
\lby{$\langle1\rangle2$, Lemma~\ref{lem:bw}, eq.~(1); no correlator of order $>2$ occurs because $\eta$ is linear
in $T(g)\Omega$ and both forms are quadratic in $\eta$}
\lstep{1}{4}{$J$ implements $r(x)=2L-x$ antiunitarily, $r'=-1$, $\{r;x\}=0$, and $T$ has weight $2$, so
$JT(x)J=T(2L-x)$ and $JT(f)J=\int \overline{f(x)}T(2L-x)dx=\int f(2L-y)T(y)dy=T(f\circ r)$ for real $f$.}
\lby{BGL93 (geometric modular conjugation); the reflection lies in the extended M\"obius group, so its
Schwarzian cocycle vanishes; $J$ antilinear and $f$ real}
\lqed{1}{5}\lby{$\langle1\rangle1$--$\langle1\rangle4$}
\begin{verdictgap}
\textbf{The mathematics of \S4.4 is correct}; the \emph{logic} as narrated is not quite what carries the weight.
\begin{itemize}[leftmargin=1.5em]
\item Linearity in $c$: valid (Prop.~\ref{prop:lin}), and there is no hidden $c$-dependence.  The Schwarzian
constant does not enter: $\widetilde T=\beta^2T$ differs from a Schwarzian-shifted field by a $c$-number, and
$(S^*-1)\Omega=0$, so $\eta$ --- hence both metrics --- is unchanged.  The choice of continuation of $g$ beyond
$L$ does not enter either: two admissible $g$'s differ by a function supported in $(L,\infty)$, so
$T(g_1)-T(g_2)$ is self-adjoint and affiliated with $\cA(I)'$, and Lemma~3.1(1) applies.  Both points should be
\emph{stated}, since a reader will ask.
\item ``\emph{The agreement of the two fixes the constant unambiguously}'' overstates what the agreement does.
Logically: (i) Prop.~\ref{prop:lin} gives $g_B=\kappa c$ for a universal $\kappa$; (ii) N7 \S5 computes
$g_B=\frac2{3\pi^2}$ for the Dirac fermion, which has $c=1$; therefore $\kappa=\frac2{3\pi^2}$.  The general
route of \S4 is \emph{not} needed for this and is not an independent determination of $\kappa$ --- it is a check
that the analytic-continuation prescription reproduces the answer.  Two computations that agree do not by
themselves ``fix'' anything unambiguously; what fixes it is (i)+(ii), and (i) is conditional on Lemma~3.2.
\item Consequently \textbf{Theorem B is conditional on Lemma~3.2 alone} (given Prop.~\ref{prop:lin} and
Theorem~A), which is a cleaner statement than the note's, and one that does not depend on the formal status of
\S4.2--4.3 at all.  I recommend restructuring \S4.4 along these lines.
\end{itemize}
\end{verdictgap}

\section{Section 5 of Note 7: the free fermion}\label{sec:fermion5}

\begin{lemma}[N7 Lemma 5.1(i), pointwise]\label{lem:point}
With $\widetilde Q(u)=-\frac i{4\pi}\sinh^{-1}\frac u2$, $u=\xi-\xi'$, and $\mathfrak g=\gt\partial_{\xi'}+\frac12\gt'$
acting on the second variable, $(Q\mathfrak g)(\xi,\xi')=\widetilde D^{(1)}(\xi,\xi')$ on $\xi<0<\xi'$.
\end{lemma}
\lstep{1}{1}{The kernel of the composition is $(Q\mathfrak g)(\xi,\xi')=-\frac12\gt'(\xi')\widetilde Q(u)-\gt(\xi')\partial_{\xi'}\widetilde Q(u)$.}
\lby{the kernel of $\partial$ is $-\partial_{\xi'}\delta(\eta-\xi')$; $\int Q(\xi,\eta)\gt(\eta)\partial_\eta\delta(\eta-\xi')d\eta
=-\partial_{\xi'}[Q(\xi,\xi')\gt(\xi')]$, and adding $\frac12\gt'Q$ leaves $-\frac12\gt'Q-\gt\partial_{\xi'}Q$.
\emph{This is where the transposition must be watched}: acting on the second variable produces $-\frac12$, not
$+\frac12$, of $\gt'$}
\lstep{1}{2}{$\gt'(\xi')=-2\sinh\xi'$, so $-\frac12\gt'\widetilde Q=\sinh\xi'\cdot\bigl(-\frac i{4\pi}\bigr)\sinh^{-1}\frac u2$;
and $\partial_{\xi'}\widetilde Q=-\partial_u\widetilde Q=-\frac i{8\pi}\frac{\cosh\frac u2}{\sinh^2\frac u2}$, so
$-\gt\partial_{\xi'}\widetilde Q=-\frac i{4\pi}\cdot\frac{2\sinh^2\frac{\xi'}2\cosh\frac u2}{\sinh^2\frac u2}$.}
\lby{$\frac d{d\xi'}(-4\sinh^2\frac{\xi'}2)=-2\sinh\xi'$; $\partial_u\sinh^{-1}\frac u2=-\frac12\frac{\cosh\frac u2}{\sinh^2\frac u2}$;
$\gt(\xi')=-4\sinh^2\frac{\xi'}2$}
\lstep{1}{3}{With $s=\frac{\xi'}2$, $U=\frac u2$:
$\frac{\sinh\xi'}{\sinh U}+\frac{2\sinh^2 s\cosh U}{\sinh^2U}=\frac{2\sinh s}{\sinh^2U}\bigl[\cosh s\sinh U+\sinh s\cosh U\bigr]
=\frac{2\sinh s\,\sinh(s+U)}{\sinh^2U}$.}
\lby{$\sinh\xi'=2\sinh s\cosh s$; addition theorem}
\lqed{1}{4}
\lby{$\langle1\rangle2$, $\langle1\rangle3$, $s+U=\frac{\xi'+\xi-\xi'}2=\frac\xi2$: the total is
$-\frac i{4\pi}\cdot\frac{2\sinh\frac{\xi'}2\sinh\frac\xi2}{\sinh^2\frac u2}
=-\frac i{2\pi}\frac{\sinh\frac\xi2\sinh\frac{\xi'}2}{\sinh^2\frac{\xi-\xi'}2}$, which is
eq.~(9) by Lemma~\ref{lem:D1}}
\begin{lemma}[N7 Lemma 5.1(ii), matrix elements]\label{lem:me}
With $\psi_\kappa(\xi)=\frac1{2\pi}e^{i\kappa\xi/2\pi}$ (the normalization forced by $\int d\kappa|\psi_\kappa\rangle\langle\psi_\kappa|=1$),
$\langle\psi_\kappa,\mathfrak g\psi_{\kappa'}\rangle=i\,\mathfrak g(\kappa,\kappa')$ with
$\mathfrak g(\kappa,\kappa')=\frac1{(2\pi)^2}\gh\bigl(\frac{\kappa-\kappa'}{2\pi}\bigr)\frac{\kappa+\kappa'}{4\pi}$,
and $D^{(1)}(\kappa,\kappa')=i(q_\kappa-q_{\kappa'})\mathfrak g(\kappa,\kappa')$.
\end{lemma}
\lstep{1}{1}{$\int d\kappa\,\psi_\kappa(\xi)\overline{\psi_\kappa(\xi')}=C^2\int d\kappa\,e^{i\kappa(\xi-\xi')/2\pi}
=C^2(2\pi)^2\delta(\xi-\xi')$, so $C=\frac1{2\pi}$.}
\lby{$\int e^{i\kappa a}d\kappa=2\pi\delta(a)$ and $\delta(y/2\pi)=2\pi\delta(y)$}
\lstep{1}{2}{$\langle\psi_\kappa,\mathfrak g\psi_{\kappa'}\rangle=\frac1{(2\pi)^2}\int_0^\infty e^{-ik\xi}\bigl[\tfrac{i\kappa'}{2\pi}\gt+\tfrac12\gt'\bigr]d\xi$
with $k=\frac{\kappa-\kappa'}{2\pi}$.}
\lby{$\langle1\rangle1$; $\mathfrak g$ is supported on $\xi>0$}
\lstep{1}{3}{$\int_0^\infty e^{-ik\xi}\gt'd\xi=\bigl[e^{-ik\xi}\gt\bigr]_0^\infty+ik\gh(k)=ik\,\gh(k)$.}
\lby{integration by parts; $\gt(0)=0$; the boundary term at $\infty$ vanishes in the regularization
$\eps>1$ of Lemma~\ref{lem:gt}, which is exactly what the ``analytic continuation'' prescription encodes ---
this is the one step of the note that is a \emph{prescription} rather than a theorem}
\lstep{1}{4}{Hence $\langle\psi_\kappa,\mathfrak g\psi_{\kappa'}\rangle=\frac1{(2\pi)^2}\frac i{2\pi}\bigl[\kappa'+\tfrac{\kappa-\kappa'}2\bigr]\gh(k)
=i\,\frac1{(2\pi)^2}\gh(k)\frac{\kappa+\kappa'}{4\pi}$.}
\lby{$\langle1\rangle2$, $\langle1\rangle3$; $\kappa'+\frac{\kappa-\kappa'}2=\frac{\kappa+\kappa'}2$}
\lstep{1}{5}{$\mathfrak g^*=-\mathfrak g$ on $(0,\infty)$, so
$\langle\psi_\kappa,(Q\mathfrak g+\mathfrak g^*Q)\psi_{\kappa'}\rangle=(q_\kappa-q_{\kappa'})\langle\psi_\kappa,\mathfrak g\psi_{\kappa'}\rangle$.}
\lby{$\langle\phi,\mathfrak g\psi\rangle=-\langle\mathfrak g\phi,\psi\rangle$ by parts (boundary terms at $0$ vanish
because $\gt(0)=0$; at $\infty$ by the same regularization); $Q\psi_\kappa=q_\kappa\psi_\kappa$}
\lstep{1}{6}{$D^{(1)}(\kappa,\kappa')$ so defined is Hermitian: $\overline{D^{(1)}(\kappa',\kappa)}=D^{(1)}(\kappa,\kappa')$
because $\overline{\gh(-k)}=\gh(k)$.}
\lby{$\gh(k)=\frac{-2i}{k(k^2+1)}$ is odd and purely imaginary}
\lqed{1}{7}\lby{$\langle1\rangle4$, $\langle1\rangle5$}
\begin{verdictok}
Lemma~5.1 verified, in both parts, with the one caveat that $\langle1\rangle3$ and $\langle1\rangle5$ discard a
boundary term at $\xi'=\infty$ under the same prescription as eq.~(7).  The note is candid about this and now
cites \texttt{rigor/kernel\_identity.tex} for an unconditional proof (digamma series), replacing the earlier claim
of a contour shift --- which I confirm does \emph{not} work: shifting $\xi'\to\xi'-i\pi$ moves the contour off the
half-line $\xi'>0$ on which $\gt$ is supported, so the shifted contour is not closed by a vanishing arc.  My own
numerics reproduce the note's four test points at $(\kappa,\kappa')=(1,3),(2.5,-1),(4,4.5),(-3,0.7)$.
\end{verdictok}

\begin{lemma}[N7 eq.~(22), the weights]\label{lem:wt}
With $u=\kappa-\kappa'$, $v=\kappa+\kappa'$, $q_\kappa=(1+e^\kappa)^{-1}$, $N=q_\kappa(1-q_{\kappa'})+q_{\kappa'}(1-q_\kappa)$:
$q_\kappa-q_{\kappa'}=-\frac{2e^{v/2}\sinh\frac u2}{(1+e^\kappa)(1+e^{\kappa'})}$,
$(1+e^\kappa)(1+e^{\kappa'})=2e^{v/2}(\cosh\frac v2+\cosh\frac u2)$,
$\frac{(q_\kappa-q_{\kappa'})^2}N=\frac{\sinh^2\frac u2}{\cosh\frac u2(\cosh\frac v2+\cosh\frac u2)}$, and
$(q_\kappa-q_{\kappa'})^2[\ell(q,q')+\ell(1-q,1-q')]=\frac{u\sinh\frac u2}{\cosh\frac v2+\cosh\frac u2}$.
\end{lemma}
\lstep{1}{1}{$q_\kappa-q_{\kappa'}=\frac{e^{\kappa'}-e^\kappa}{(1+e^\kappa)(1+e^{\kappa'})}$ and
$e^{\kappa'}-e^\kappa=e^{v/2}(e^{-u/2}-e^{u/2})=-2e^{v/2}\sinh\frac u2$.}
\lby{$\kappa=\frac{v+u}2$, $\kappa'=\frac{v-u}2$; \emph{the minus sign is the one the 16:13 erratum inserted, and
it is correct}}
\lstep{1}{2}{$\cosh\frac\kappa2\cosh\frac{\kappa'}2=\frac12(\cosh\frac v2+\cosh\frac u2)$ and
$(1+e^\kappa)(1+e^{\kappa'})=4e^{v/2}\cosh\frac\kappa2\cosh\frac{\kappa'}2$.}
\lby{product-to-sum; $1+e^\kappa=2e^{\kappa/2}\cosh\frac\kappa2$}
\lstep{1}{3}{$N=\frac{e^\kappa+e^{\kappa'}}{(1+e^\kappa)(1+e^{\kappa'})}$ and $e^\kappa+e^{\kappa'}=2e^{v/2}\cosh\frac u2$.}
\lby{$1-q_\kappa=\frac{e^\kappa}{1+e^\kappa}$}
\lstep{1}{4}{$\frac{(q-q')^2}N=\frac{4e^v\sinh^2\frac u2}{\mathcal D^2}\cdot\frac{\mathcal D}{2e^{v/2}\cosh\frac u2}
=\frac{2e^{v/2}\sinh^2\frac u2}{\cosh\frac u2\cdot2e^{v/2}(\cosh\frac v2+\cosh\frac u2)}$, $\mathcal D:=(1+e^\kappa)(1+e^{\kappa'})$.}
\lby{$\langle1\rangle1$--$\langle1\rangle3$}
\lstep{1}{5}{$(q-q')^2[\ell(q,q')+\ell(1-q,1-q')]=(q_\kappa-q_{\kappa'})\bigl[\log\frac{q_\kappa}{1-q_\kappa}-\log\frac{q_{\kappa'}}{1-q_{\kappa'}}\bigr]
=(q_\kappa-q_{\kappa'})(\kappa'-\kappa)$.}
\lby{$\ell(a,b)=\frac{\log a-\log b}{a-b}$, $(1-q_\kappa)-(1-q_{\kappa'})=-(q_\kappa-q_{\kappa'})$, and
$\log\frac{q}{1-q}=-\kappa$}
\lqed{1}{6}
\lby{$\langle1\rangle5$ with $\langle1\rangle1$: $(-\frac{2e^{v/2}\sinh\frac u2}{\mathcal D})(-u)
=\frac{u\sinh\frac u2}{\cosh\frac v2+\cosh\frac u2}$, manifestly $\ge0$ for both signs of $u$}
\begin{verdictok}
Eq.~(22) verified, both lines, including the identity $\cosh\frac\kappa2\cosh\frac{\kappa'}2=\frac12(\cosh\frac v2+\cosh\frac u2)$
and the sign corrected by the erratum.  Also verified: $|\mathfrak g|^2=\frac{|\gh(u/2\pi)|^2v^2}{256\pi^6}$
(from $\mathfrak g=\frac1{4\pi^2}\gh\cdot\frac v{4\pi}$, $16\pi^4\cdot16\pi^2=256\pi^6$),
$|\gh(\frac u{2\pi})|^2=\frac{256\pi^6}{u^2(u^2+4\pi^2)^2}$ (from $|\gh(k)|^2=\frac4{k^2(k^2+1)^2}$ at $k=\frac u{2\pi}$),
and $d\kappa\,d\kappa'=\frac12du\,dv$.
\end{verdictok}
\begin{lemma}[N7 eq.~(23): the base integral, proved]\label{lem:base}
For real $a\ne0$, $\displaystyle\int_{\mathbb R}\frac{w^2\,dw}{\cosh w+\cosh a}=\frac{2a(a^2+\pi^2)}{3\sinh a}$,
and consequently $V(u)=\int_{\mathbb R}\frac{v^2dv}{\cosh\frac v2+\cosh\frac u2}=\frac{2u(u^2+4\pi^2)}{3\sinh\frac u2}$.
\end{lemma}
\lstep{1}{1}{$F(z)=\frac{z^3}{\cosh z+\cosh a}$ has, in the strip $0<\operatorname{Im}z<2\pi$, exactly the two
simple poles $z_\pm=\pm a+i\pi$, with $\cosh'(z_\pm)=\sinh(\pm a+i\pi)=\mp\sinh a$.}
\lby{$\cosh z=-\cosh a=\cosh(a+i\pi)\Leftrightarrow z\in\{\pm(a+i\pi)\}+2\pi i\mathbb Z$; simplicity because $\sinh a\ne0$}
\lstep{1}{2}{Over the rectangle $[-R,R]\times[0,2\pi]$, the two horizontal edges give
$\int_{\mathbb R}\frac{w^3-(w+2\pi i)^3}{\cosh w+\cosh a}dw=-6\pi i\,I+12\pi^2\!\!\int\!\frac{w\,dw}{\cosh w+\cosh a}+8\pi^3i\,I_0
=-6\pi i\,I+8\pi^3i\,I_0$, where $I=\int\frac{w^2dw}{\cosh w+\cosh a}$, $I_0=\int\frac{dw}{\cosh w+\cosh a}$.}
\lby{$\cosh(z+2\pi i)=\cosh z$; binomial expansion; the $w$-integral vanishes by oddness}
\lstep{1}{3}{The vertical edges vanish as $R\to\infty$.}
\lby{$|\cosh(\pm R+iy)|\ge\sinh R$, numerator $O(R^3)$, edge length $2\pi$}
\lstep{1}{4}{$2\pi i\bigl[\operatorname{Res}_{z_+}+\operatorname{Res}_{z_-}\bigr]
=2\pi i\frac{-(a+i\pi)^3+(-a+i\pi)^3}{\sinh a}=2\pi i\frac{-2a^3+6\pi^2a}{\sinh a}$.}
\lby{$\langle1\rangle1$; $(a+i\pi)^3=a^3+3i\pi a^2-3\pi^2a-i\pi^3$ and $(-a+i\pi)^3=-a^3+3i\pi a^2+3\pi^2a-i\pi^3$}
\lstep{1}{5}{The same argument with $F(z)=\frac{z}{\cosh z+\cosh a}$ gives $-2\pi i\,I_0=2\pi i\frac{-2a}{\sinh a}$,
i.e.\ $I_0=\frac{2a}{\sinh a}$.}
\lby{$\langle1\rangle2$--$\langle1\rangle4$ with $z^3$ replaced by $z$}
\lstep{1}{6}{$-6\pi I+8\pi^3I_0=2\pi\frac{-2a^3+6\pi^2a}{\sinh a}$, hence
$-6\pi I=\frac{-4\pi a^3+12\pi^3a-16\pi^3a}{\sinh a}=\frac{-4\pi a(a^2+\pi^2)}{\sinh a}$, i.e.\
$I=\frac{2a(a^2+\pi^2)}{3\sinh a}$.}
\lby{$\langle1\rangle2$, $\langle1\rangle4$, $\langle1\rangle5$; divide by $i$}
\lqed{1}{7}
\lby{$\langle1\rangle6$; substituting $v=2w$ in $V(u)$ gives $8\int\frac{w^2dw}{\cosh w+\cosh\frac u2}
=8\cdot\frac{2\cdot\frac u2(\frac{u^2}4+\pi^2)}{3\sinh\frac u2}=\frac{2u(u^2+4\pi^2)}{3\sinh\frac u2}$}
\begin{verdictok}
The note calls eq.~(23) ``\emph{verified numerically}''; it is now proved.  Numerical confirmation:
$a=0.5$: $6.4732908384842$ both sides; $a=2$: $5.0988468781728$ both sides;
$V(0.7)=52.218545714491$, $V(3)=45.53504471981$, both matching the closed form to 14 digits.
\end{verdictok}

\begin{theorem}[N7 Theorem 5.2]\label{thm:A5}
For the $r$-component free chiral fermion, $g_B=\frac23\int_{\mathbb R}\frac{\tanh\frac u2}{u(u^2+4\pi^2)}du\cdot r=\frac{2r}{3\pi^2}$
and $g_{\mathrm{KM}}=\frac13\int_{\mathbb R}\frac{du}{u^2+4\pi^2}\cdot r=\frac r6$.
\end{theorem}
\lstep{1}{1}{$g_B=2\iint d\kappa d\kappa'\frac{|D^{(1)}|^2}N=\iint du\,dv\,|\mathfrak g|^2\frac{(q_\kappa-q_{\kappa'})^2}N
=\frac1{256\pi^6}\int du\,|\gh(\tfrac u{2\pi})|^2\frac{\sinh^2\frac u2}{\cosh\frac u2}V(u)$.}
\lby{eq.~(20), Lemma~\ref{lem:me}, Lemma~\ref{lem:wt}, $d\kappa d\kappa'=\frac12dudv$ ($2\cdot\frac12=1$), and the
$v$-integral of Lemma~\ref{lem:base}}
\lstep{1}{2}{$\frac1{256\pi^6}\cdot\frac{256\pi^6}{u^2(u^2+4\pi^2)^2}\cdot\frac{\sinh^2\frac u2}{\cosh\frac u2}
\cdot\frac{2u(u^2+4\pi^2)}{3\sinh\frac u2}=\frac23\frac{\tanh\frac u2}{u(u^2+4\pi^2)}$.}
\lby{$\langle1\rangle1$, Lemma~\ref{lem:wt}, Lemma~\ref{lem:base}: all powers of $\pi$ cancel, as the note says}
\lstep{1}{3}{With $u=2\pi k$: $\int_{\mathbb R}\frac{\tanh\frac u2}{u(u^2+4\pi^2)}du
=\int\frac{2\pi\,dk\,\tanh\pi k}{2\pi k\cdot8\pi^2(k^2+1)}\cdot\frac{8\pi^2}{8\pi^2}
=\frac1{4\pi^2}\int_{\mathbb R}\frac{\tanh\pi k}{k(k^2+1)}dk=\frac1{4\pi^2}\cdot4=\frac1{\pi^2}$.}
\lby{$u(u^2+4\pi^2)=2\pi k\cdot4\pi^2(k^2+1)=8\pi^3k(k^2+1)$, $du=2\pi dk$, so the prefactor is
$\frac{2\pi}{8\pi^3}=\frac1{4\pi^2}$; and Lemma~\ref{lem:ints}$\langle1\rangle6$ gives $\int_{\mathbb R}=4$}
\lstep{1}{4}{$g_{\mathrm{KM}}=\iint d\kappa d\kappa'|\mathfrak g|^2(q_\kappa-q_{\kappa'})^2[\ell+\ell]
=\frac1{512\pi^6}\int du\,|\gh|^2u\sinh\tfrac u2\,V(u)=\frac13\int\frac{du}{u^2+4\pi^2}=\frac13\cdot\frac12=\frac16$.}
\lby{Lemma~\ref{lem:wt} second line, $\frac12$ from the Jacobian, Lemma~\ref{lem:base}; the $\sinh\frac u2$
cancels against $V(u)$'s $1/\sinh\frac u2$, leaving $\frac1{512\pi^6}\cdot256\pi^6\cdot\frac23\frac{u^2}{u^2(u^2+4\pi^2)}$;
$\int_{\mathbb R}\frac{du}{u^2+4\pi^2}=\frac\pi{2\pi}=\frac12$}
\lqed{1}{5}\lby{$\langle1\rangle2$--$\langle1\rangle4$; $g_B=\frac23\cdot\frac1{\pi^2}=\frac2{3\pi^2}$; components add}
\begin{verdictok}
Theorem~5.2 verified.  Numerically $\int_{\mathbb R}\frac{\tanh(u/2)}{u(u^2+4\pi^2)}du=0.101321183642338=1/\pi^2$
to 15 digits, $g_B=0.0675474557615585=2/(3\pi^2)$, $\int_{\mathbb R}\frac{du}{u^2+4\pi^2}=0.5$,
$g_{\mathrm{KM}}=0.1666666666\ldots$.  \textbf{I confirm the erratum}: the proof text must read
$\frac1{4\pi^2}\int_{\mathbb R}\frac{\tanh\pi k}{k(k^2+1)}dk=\frac1{\pi^2}$, and the version now on disk does.
(The earlier ``$\frac1{\pi^2}\int\ldots=\frac4{\pi^2}$'' was wrong by a factor $4$ in \emph{both} the prefactor
and the value, so the product --- and the boxed result --- was always right.)  Note also that
Theorem~5.2's $g_B$ is \emph{exactly} $c=1$ of eq.~(17): $\frac1{6\pi^2}\int_{\mathbb R}\frac{\tanh\pi k}{k(k^2+1)}dk
=\frac23\cdot\frac1{4\pi^2}\int_{\mathbb R}(\cdots)$, the same integral --- so the two routes really do coincide
integrand by integrand, not merely in the final number.
\end{verdictok}
\section{Section 7 of Note 7: Corollary 6.1}\label{sec:cor}
\begin{verdictok}
The conversions are correct.  Per chirality, $\zeta_0=z(1+e^0)=2z$ and $-\log\Fid=\frac c{12\pi^2}(2z)^2=\frac c{3\pi^2}z^2$;
with $z=\eta_{\mathrm V}/(1-\eta_{\mathrm V})=\eta_{\mathrm V}+O(\eta_{\mathrm V}^2)$ the two chiralities give
$\frac{c+\bar c}{3\pi^2}\eta^2$, numerically $\frac2{3\pi^2}c=0.067547\,c$ (the note's $0.06755$);
$D=\frac c{12}(2z)^2\cdot2=\frac{c+\bar c}3\eta^2=0.6667\,c\eta^2$.  With $I=\frac c6\log(1+z)=\frac c6z+O(z^2)$,
i.e.\ $z=6I/c+O(I^2)$, one gets $-\log\Fid=\frac c{3\pi^2}\frac{36I^2}{c^2}=\frac{12}{\pi^2c}I^2$.  The
two-branch law $-\log F^{(\lambda)}=\Phi_c(z(1+e^{-\pi\lambda}))+\Phi_{\bar c}(z(1+e^{\pi\lambda}))$ is even in
$\lambda$ for $c=\bar c$ and, since $\frac d{d\lambda}[(1+e^{-\pi\lambda})^2+(1+e^{\pi\lambda})^2]
=2\pi[e^{\pi\lambda}+e^{2\pi\lambda}-e^{-\pi\lambda}-e^{-2\pi\lambda}]>0$ for $\lambda>0$, minimal at $\lambda=0$;
so ``the ordinary Petz map is the optimum of the family'' is right at quadratic order.  The chiral factor-$4$
statement ($\zeta_t\downarrow z$) is right and now correctly labelled chiral-only.
\end{verdictok}
\begin{verdictgap}
Two presentational risks.  (i) ``$-\log\Fid\simeq\frac{12}{\pi^2c}I^2$'' sits immediately after the
two-dimensional discussion but is the \emph{chiral} law ($I=\frac c6\log(1+z)$ is the chiral CMI); for a 2D CFT
with $c=\bar c$ and $I_{\mathrm{tot}}=\frac c3\log(1+z)$ the same computation gives
$-\log F=\frac6{\pi^2c}I_{\mathrm{tot}}^2$, i.e.\ $\frac{12}{\pi^2c_{\mathrm{tot}}}I_{\mathrm{tot}}^2$ with
$c_{\mathrm{tot}}=c+\bar c$.  One sentence would remove the ambiguity.  (ii) The corollary is stated ``for every
diffeomorphism covariant chiral net'' but is conditional on Lemma~3.2 (Theorem~B) and, for the free fermion, on
Note~5 Cor.~4.2 (the interchange of the compression limit with the second-order expansion) --- the main box now
says this for Theorem~A, but Corollary~6.1 itself only says ``(Theorem~B; rigorous for the free fermion with
$c=r$)'', which overstates the free-fermion case: what is rigorous there is the pair of \emph{metrics}, not yet
$\lim\Phi/\zeta^2$.
\end{verdictgap}

\section{Note 6: the mathematical and citation claims}\label{sec:n6}

\subsection{\S2 (P1): an internal normalization clash}
\begin{issue}[N6 \S2, step 3 vs step 4]
Step~3 states ``$\langle\widetilde T(\xi)\widetilde T(\xi')\rangle=\frac c{32}\sinh^{-4}\frac{\xi-\xi'-i0}2$'';
step~4 (as amended) states ``\emph{With the generator normalization $\langle T(x)T(y)\rangle=\frac c{8\pi^2}(x-y-i0)^{-4}$
(Note~7, eq.~(1))}''.
\end{issue}
\lstep{1}{1}{By Lemma~\ref{lem:W}, the generator normalization gives
$\langle\widetilde T\widetilde T\rangle=\frac c{128\pi^2}\sinh^{-4}\frac{u-i0}2$.}
\lby{Lemma~\ref{lem:W}}
\lstep{1}{2}{$\frac c{32}\big/\frac c{128\pi^2}=4\pi^2$, which is the ratio $\bigl(\frac{c/2}{c/8\pi^2}\bigr)$
between the BPZ and the generator two-point constants.}
\lby{arithmetic; $(2\pi)^2=4\pi^2$}
\lqed{1}{3}\lby{$\langle1\rangle1$, $\langle1\rangle2$: step~3 uses $\langle TT\rangle=\frac c2(x-y)^{-4}$}
\begin{verdictbreak}
\textbf{Confirms [V6].}  The 16:13 erratum fixed step~4 but left step~3 at the BPZ constant $\frac c{32}$, so the
two steps of the \emph{same} argument now use two normalizations differing by $4\pi^2$.  Nothing downstream in
Note~6 is numerically wrong --- eq.~(4) is stated with ``$\propto$'' and the final numbers are imported from
Note~7 --- but as it stands \S2 does not derive what it claims to derive, and a reader following step~3 into
step~4 will land on $f_2=4\pi^2\cdot\frac c{12\pi^2}=\frac c3$.  Repair: $\frac c{32}\to\frac c{128\pi^2}$, and
correspondingly ``the weight-2 thermal spectral density $\propto c\,k(k^2+1)/(1-e^{-2\pi k})$'' should be given
with its constant, $\Wh(k)=\frac c{24\pi}\frac{k(k^2+1)}{e^{2\pi k}-1}$ (Lemma~\ref{lem:What7}).  Also: the
sentence in step~4 ``\emph{This step has so far been fixed by matching the free-fermion number}'' is now stale ---
Note~7 derives it; and ``\emph{deriving it directly from the spectral density of $\widetilde T$ is the first task
below}'' should be marked done.
\end{verdictbreak}
\begin{verdictok}
The rest of N6 \S2 checks out.  Eq.~(2): $g_B=2\sum|A_{ij}|^2\frac{(p_i-p_j)^2}{p_i+p_j}=2\langle A\Omega,\frac{(\Delta-1)^2}{\Delta+1}A\Omega\rangle$
and $g_{\mathrm{KM}}=\sum|A_{ij}|^2(p_i-p_j)\log\frac{p_i}{p_j}=\langle A\Omega,(\Delta-1)\log\Delta\,A\Omega\rangle$
are both correct (put $p_jF(\lambda)$ equal to the summand as in Lemma~\ref{lem:std}$\langle1\rangle3$);
they are the $A\in\cM$ specializations of N7 eq.~(13), and carry the same caveat.
Eq.~(4): $\int_0^\infty\frac{\tanh(\kappa/2)}{\kappa(\kappa^2+4\pi^2)}d\kappa=\frac1{4\pi^2}\int_0^\infty\frac{\tanh\pi k}{k(k^2+1)}dk=\frac2{4\pi^2}=\frac1{2\pi^2}$,
so $\int_{\mathbb R}=\frac1{\pi^2}$ as stated; $\tanh\frac\kappa2=\sum_{n\ge0}\frac{4\kappa}{\kappa^2+(2n+1)^2\pi^2}$
is the Mittag-Leffler series at $z=\kappa/2$; $\sum_{n\ge0}\frac1{(2n+1)(2n+3)}=\frac12$; and
$\int_{\mathbb R}\frac{d\kappa}{\kappa^2+4\pi^2}=\frac12$.  The ``ratio of the two integrals \dots\ is exactly
$\pi^2$, independent of normalizations'' is correct.
\end{verdictok}
\subsection{\S3 (P2): Observation 1, the energy formula, Route B}

\begin{verdictok}[Observation 1, first arrow]
``$k_s$ is $C^{1,1}$ at the touching point'' is right: $h_s(0)=0$, $h_s'(x)=\frac{L^2}{(L+sx)^2}$ so $h_s'(0)=1$,
$h_s''(0)=-2s/L\ne0=\mathrm{id}''(0)$, so the glued map has a jump of the second derivative and no worse.
``$\log\widetilde k'$ Lipschitz, hence in $H^1(S^1)\subset H^{1/2}(S^1)$'' is right: $\widetilde k'$ is Lipschitz
and bounded away from $0$ on the compact $S^1$, so $\log\widetilde k'$ is Lipschitz, its distributional derivative
is in $L^\infty\subset L^2(S^1)$, and $H^1(S^1)\hookrightarrow H^{1/2}(S^1)$ by the Fourier characterization.
Existence of the extension $\widetilde k$ is also right (prescribe the $1$-jets at both endpoints of $I'$ and
interpolate); it is far from unique, but $\operatorname{Ad}U(\widetilde k)|_{\cA(I)}$ is independent of the choice
(two extensions differ by a diffeomorphism trivial on $I$, whose implementer lies in $\cA(I)'$), while the
\emph{vector} $U(\widetilde k)^*\Omega$ is not --- which is why Conjecture~3.2's ``optimized over the extension''
is the right formulation.
\end{verdictok}
\begin{verdictgap}[Observation 1, the implementability arrow]
I confirm [R4]: ``WP $\Rightarrow$ implementable in the fermionic Fock space'' is not supplied by the cited
sources, and the amended text now says so.  Two additions.  (i) The theorem that is actually wanted is
Shale--Stinespring: a Bogoliubov transformation of the CAR algebra is unitarily implementable in the Fock
representation of $P_+$ iff $[P_+,\Gamma]$ is Hilbert--Schmidt; the missing link is that the \emph{weight-$\frac12$}
action $\phi\mapsto(\phi\circ\varphi^{-1})\sqrt{(\varphi^{-1})'}$ has HS off-diagonal blocks iff
$\log\varphi'\in H^{1/2}$.  (ii) The sentence ``\emph{So $U(\widetilde k)$ exists in the fermionic Fock space}''
still begins with ``So'', i.e.\ still presents as a deduction what the preceding sentence has just declared
missing; it should read ``Granting this, $U(\widetilde k)$ exists\ldots''.
\end{verdictgap}

\begin{lemma}[N6 eq.~(5)]\label{lem:energy}
Let $L_0=\int_0^{2\pi}\mathcal T(\theta)d\theta$ be the conformal Hamiltonian ($L_0\Omega=0$, $\langle\mathcal T\rangle=0$).
Then, for $\varphi\in\Diff_+(S^1)$ and $\psi:=\varphi^{-1}$,
\[
 E(\varphi)=\langle U(\varphi)\Omega,L_0U(\varphi)\Omega\rangle=\frac c{48\pi}\int_0^{2\pi}\bigl(\partial_\theta\log\psi'\bigr)^2d\theta
 =\frac c{48\pi}\int_0^{2\pi}\Bigl[\bigl(\partial_\theta\log\psi'\bigr)^2-\bigl(\psi'^2-1\bigr)\Bigr]d\theta
 +\frac c{48\pi}\int_0^{2\pi}(\psi'^2-1)d\theta .
\]
\end{lemma}
\lstep{1}{1}{$U(\varphi)^*\mathcal T(\theta)U(\varphi)=\bigl(\psi'(\theta)\bigr)^2\mathcal T(\psi(\theta))+A\{\psi;\theta\}$
with $A=-\frac c{24\pi}$.}
\lby{weight-2 covariance (Lemma~\ref{lem:bw}$\langle1\rangle4$) plus the Virasoro cocycle; $A$ is fixed by the
$c$-number part of $[\mathcal T(\theta'),\mathcal T(\theta)]$, which by eq.~(1) is
$\frac c{8\pi^2}\bigl[(u-i0)^{-4}-(u+i0)^{-4}\bigr]=-\frac{ic}{24\pi}\delta'''(\theta'-\theta)$, giving
$\langle[T(h),\mathcal T(\theta)]\rangle=\frac{ic}{24\pi}h'''(\theta)=-iAh'''(\theta)$}
\lstep{1}{2}{$\langle\mathcal T\rangle=0$ by rotation invariance and $L_0\Omega=0$, so
$E(\varphi)=A\int_0^{2\pi}\{\psi;\theta\}d\theta$.}
\lby{$\langle1\rangle1$; the first term has vanishing vacuum expectation}
\lstep{1}{3}{With $\Lambda=\log\psi'$: $\{\psi;\theta\}=\Lambda''-\frac12(\Lambda')^2$, and
$\int_0^{2\pi}\Lambda''=0$, so $\int\{\psi;\theta\}d\theta=-\frac12\int(\Lambda')^2d\theta$.}
\lby{$\frac{\psi'''}{\psi'}=\Lambda''+(\Lambda')^2$ and $\frac32\bigl(\frac{\psi''}{\psi'}\bigr)^2=\frac32(\Lambda')^2$;
periodicity}
\lqed{1}{4}\lby{$\langle1\rangle2$, $\langle1\rangle3$: $E=-\frac c{24\pi}\cdot(-\frac12)\int(\Lambda')^2=\frac c{48\pi}\int(\Lambda')^2$}
\begin{verdictgap}
Two findings.
(i) \textbf{The displayed combination is correct} in the convention in which $L_0$ is the CFT generator with
$\langle T_{\mathrm{cyl}}\rangle=-\frac c{24}$ per $2\pi$, i.e.\ eq.~(5) computes the energy \emph{relative to the
Casimir value}; equivalently, with the AQFT $L_0$ ($L_0\Omega=0$) the two forms agree because
$-\frac c{48\pi}\int(\psi'^2-1)d\theta$ is exactly the difference of the two conventions.  Numerically I confirm
that the combination vanishes identically on $\mathrm{PSL}(2,\mathbb R)$: for the M\"obius circle maps
$\tan\frac\varphi2=\lambda\tan\frac\theta2$, $\lambda=1.7$ and $\lambda=1/1.7$,
$\int(\partial\log\varphi')^2=\int(\varphi'^2-1)=0.905517882505$, difference $<3\times10^{-21}$.  So the note's
hedge ``\emph{the sign of the second term depends on conventions}'' can be replaced by the assertion that the
displayed sign is the right one.
(ii) \textbf{The argument is $\varphi^{-1}$, not $\varphi$.}  $E(\varphi)$ involves $\log(\varphi^{-1})'$, and the
two are genuinely different: for $\varphi(\theta)=\theta+0.21\sin2\theta$ the bracket integrates to
$2.006827$ with $\varphi$ and $2.325495$ with $\varphi^{-1}$.  Since Route~B applies the bound to
$U(\varphi_n^{-1})\Omega$, the practical consequence is nil, but eq.~(5) as displayed is $E(\varphi^{-1})$.
Repair: write $\psi=\varphi^{-1}$ in eq.~(5), or state $E(\varphi)=\frac c{48\pi}\int[(\partial\log\varphi')^2-(\varphi'^2-1)]d\theta$
for $\langle U(\varphi)^*\Omega,L_0U(\varphi)^*\Omega\rangle$.
(iii) ``\emph{It is finite iff $\log\varphi'\in H^1$}'' is right for the first term; the second needs only
$\varphi'\in L^2$, automatic for a $C^1$ diffeomorphism.  ``$C^{1,1}$'' suffices, as claimed.
\end{verdictgap}
\begin{verdictgap}[Route B: the compactness step]
``\emph{the set $\{\psi:\langle\psi,L_0\psi\rangle\le E\}$ is norm-compact}'' is \textbf{false as written}: the
set contains $\mathbb C\Omega$ (since $L_0\Omega=0$) and is therefore unbounded.  What is true, and what the
argument uses, is that $\{\psi:\|\psi\|\le1,\ \langle\psi,L_0\psi\rangle\le E\}$ is norm-compact when $L_0$ has
compact resolvent: with $P_N$ the spectral projection of $L_0$ onto $[0,N]$, $\|(1-P_N)\psi\|^2\le E/N$ uniformly
on the set, so it is totally bounded, and it is norm-closed because $\psi\mapsto\langle\psi,L_0\psi\rangle$ is
norm-lower-semicontinuous.  The vectors $U(\varphi_n^{-1})\Omega$ are unit vectors, so the repair is one clause.
Also, ``finite-dimensional $L_0$-eigenspaces'' alone does not give compact resolvent; one needs in addition that
the eigenvalues accumulate only at $+\infty$, which does hold for the vacuum module of a unitary VOA of CFT type
($\mathrm{spec}\,L_0\subset\mathbb Z_{\ge0}$ with finite multiplicities).
\end{verdictgap}
\begin{verdictgap}[Route B: the identification step]
``\emph{the implementers of diffeomorphisms that are trivial on a region commute with its algebra}'' is correct
but is a \emph{theorem}, not a definition, and is uncited: for a diffeomorphism covariant net,
$\operatorname{supp}\varphi\subset K$ implies $U(\varphi)\in\cA(K)$; hence $\varphi|_J=\mathrm{id}$ gives
$U(\varphi)\in\cA(J')\subset\cA(J)'$.  Locator to add: Carpi--Weiner, CMP \textbf{258} (2005), Prop.~3.7 (or
D'Antoni--Fredenhagen--K\"oster, Lett.\ Math.\ Phys.\ \textbf{67} (2004)).  Two further hypotheses are used but
not carried in Conjecture~3.2: \emph{strong additivity} (to make $\cA(A)\vee\cA(D)=\cA(I)$ and hence the
$*$-algebra $\bigcup_{A'\Subset A,D'\Subset D}\cA(A')\vee\cA(D')$ weakly dense in $\cA(I)$ --- the note invokes it
in the text but the conjecture says only ``finite $L_0$-multiplicities''), and \emph{Haag duality}.  The final
inference is sound: $\omega\circ\beta_s$ is only a positive functional on the dense $*$-algebra, but it agrees
there with the normal state $\omega_\Psi$, so it \emph{is} $\omega_\Psi$ and is normal.  The step
``independence of the subsequence'' is listed as to-be-done in the programme, correctly.
\end{verdictgap}

\subsection{\S4.1 (P3) and the CMI-squared coefficient}
\begin{verdictok}
The conversion of Vardhan--Wei--Zou's $-\log F^{(0)}=-\frac c9\log(1-\eta)+O(1)$ into
$\Phi(\zeta)\simeq\frac c{18}\log\zeta$ per chirality is correct and, importantly, remains correct under the
corrected two-branch law: at $\lambda=0$ both chiralities carry the same $\zeta_\pm=2z$, so
$-\log F^{(0)}=2\Phi_c(2z)=\frac c9\log\zeta+O(1)$, matching $-\frac c9\log(1-\eta)=\frac c9\log(1+z)$ with
$1-\eta=(1+z)^{-1}$ and $\zeta=2z$.  The coefficient $\frac{12}{\pi^2c}$ of $I^2$ in N6 \S1 and Conjecture~2.1 is
the chiral one, consistent with N7 Cor.~6.1 and with $z=6I/c+O(I^2)$; for a 2D CFT the same statement reads
$\frac{12}{\pi^2(c+\bar c)}I_{\mathrm{tot}}^2$.
\end{verdictok}
\begin{verdictgap}
Conjecture~2.1 still asserts ``$-\log\Fid=\frac c{12\pi^2}\zeta_t^2[1+O(\log^{-2}(1/\zeta_t))]$'' and ``for a
two-dimensional CFT the two chiralities add'', without the opposite-rotation qualification that the 16:13 errata
put into N5 Cor.~5.1 and N7 Cor.~6.1; and its ``$\zeta=2z$'' is the $\lambda=0$ value only.  For consistency
across the series, Conjecture~2.1 should carry the two-branch form
$-\log F^{(\lambda)}=\Phi_c(z(1+e^{-\pi\lambda}))+\Phi_{\bar c}(z(1+e^{\pi\lambda}))$.
\end{verdictgap}

\section{Summary of findings}\label{sec:summary}
\begin{enumerate}[leftmargin=1.6em]
\item[\textbf{BREAK}] \textbf{N7, sentence after eq.~(13)}: ``both are bounded by $1+\lambda$ \dots\ the second
variations are finite whenever $A\Omega$ is a vector''.  False for the Kubo--Mori weight
($(\lambda-1)\log\lambda/(1+\lambda)\to\infty$).  Local; no result falls.  (\S\ref{sec:bound})
\item[\textbf{BREAK}] \textbf{N6 \S2 step 3}: ``$\langle\widetilde T\widetilde T\rangle=\frac c{32}\sinh^{-4}$''
is the BPZ constant, $4\pi^2$ times the generator-normalized value $\frac c{128\pi^2}$ that step~4 now invokes.
The two halves of one argument use two normalizations.  (\S\ref{sec:n6})
\item[\textbf{GAP}] \textbf{N7 \S4.2}: the KMS sign-fixing presupposes $T(g)\in\cM$, which is false; the sign is
nevertheless the note's, by Bisognano--Wichmann (Lemma~\ref{lem:bw}).  (\S\ref{sec:kmssign})
\item[\textbf{GAP}] \textbf{N7 \S1}: ``a translation $\xi\mapsto\xi\pm2\pi t$'' --- the sign must be fixed, and
it is $-$.  (\S\ref{sec:bwsign})
\item[\textbf{GAP}] \textbf{N7 \S4.2--4.3 as a whole}: eq.~(13) is the $A\in\cM$ formula and no admissible
representative is in $\cM$; the general route is formal and is validated only by \S5.  (\S\ref{sec:normdiv})
\item[\textbf{GAP}] \textbf{N7 Lemma 3.2}: ``affiliated with a von Neumann algebra containing $\cM$'' is a
vacuous hypothesis; ``$o(s)$ controlled in the norm of the tangent space'' is undefined; the form-domain
condition is needed only for $g_{\mathrm{KM}}$.
\item[\textbf{GAP}] \textbf{N7 \S4.4}: ``the agreement of the two fixes the constant unambiguously'' --- what
fixes it is linearity in $c$ (Prop.~\ref{prop:lin}) plus Theorem~5.2; \S4 is a check, not a determination.
\item[\textbf{GAP}] \textbf{N6 \S3 Route B}: the ``norm-compact'' set is unbounded as written; two uncited
theorems (implementers of localized diffeomorphisms; compact resolvent from VOA gradation); strong additivity
missing from Conjecture~3.2's hypotheses.
\item[\textbf{MINOR}] \textbf{N6 eq.~(5)}: the displayed functional is $E(\varphi^{-1})$, not $E(\varphi)$; the
relative sign of the two terms is however correct and the combination does vanish on $\mathrm{PSL}(2,\mathbb R)$.
\item[\textbf{MINOR}] N7 after eq.~(6): $\langle\widetilde T\rangle=0$ in the note's own convention, so
``$\widetilde T_0=\widetilde T-\langle\widetilde T\rangle$'' is vacuous; N7 Lemma~5.1(i): ``on $A\times A$ and
$D\times D$ both sides vanish'' needs the M\"obius-invariance argument; N7 Cor.~6.1: the $\frac{12}{\pi^2c}$ law
is chiral; N6 Conjecture~2.1 lacks the two-branch qualification; N6 Observation~1: ``So $U(\widetilde k)$
exists'' still reads as a deduction from what was just declared missing.
\item[\textbf{OK}] Everything else, in particular: eq.~(1) (proved independently, Theorem~\ref{thm:TT}),
eqs.~(2),(3),(6)--(9), Lemma~3.1 in full, Lemma~4.1, eqs.~(16)--(19), Lemma~5.1, eqs.~(22),(23), Theorem~5.2 and
its corrected $\frac1{\pi^2}$, Cor.~6.1's conversions, N6 eqs.~(2),(4) and the $c/18$ conversion.
\end{enumerate}
\end{document}
